\chapter{用户与对象理解}
\label{chap:rs-understanding}

某位用户昨天浏览了登山包、水壶和相机，今天要求“防水通勤背包，价格不超过500元”。系统已经有了一些记录，却还没有得到可以直接推荐的答案：昨天的相机浏览是否与今天有关，商品标题中的“户外”是否意味着适合通勤，未注明重量的背包是否满足轻便要求，都需要分别判断。本章把这些判断所需的信息组织起来，说明怎样从已知证据形成需求、对象及两者匹配的表示。

读者已经熟悉向量表示、监督学习以及循环网络、注意力和图网络。本章关心这些工具在推荐中接收什么证据、用什么监督学习、最终输出什么。相同历史可以被压成一个状态、保留为多个兴趣，也可以等候选到来后再选择相关部分；三种做法承担的计算接口不同。信息不足时，还需决定能否补充内容、迁移已有知识或取得新反馈。只有把这些条件说清楚，后续召回和排序才能知道自己实际拿到了什么依据。

% Outline: 2.1
\section{理解任务与信息来源}
\label{sec:rs-understanding-evidence}

理解环节提供三类相互关联的输出。用户侧给出需求与偏好依据，例如当前用途、价格上限和近期摄影行为；对象侧给出属性与关系，例如容量、重量、接口兼容性及供货状态；匹配侧把需求对应到对象，例如“价格合格、重量未知、通勤适用性有待核验”。匹配结果可以是结构化条件，也可以是服务某项预测的向量或分数。向量中的某一维没有被定义为重量时，就不能将其数值当作重量检查结果。

证据需要保留来源。用户亲自填写“不超过500元”是直接表达；日志记录了相机点击，是交互观测；模型由此推测“可能需要摄影用品”，则是推断。直接表达也可能含混或过时，但不能仅因历史模型给高价商品打了高分就被悄悄覆盖。交互观测还受到机会影响：用户点击了看到的商品，却未必有机会比较未展示的商品。因此，训练一个行为预测器与证明某项偏好成立之间仍有距离。

用户之间的关注、好友和信任关系可以提供额外依据，但它们与交互相似性不是同一关系。两人互相关注，不意味着购买过相同物品；两人买过相同相机，也不证明存在社交联系。对象之间的品牌、部件和主题关系又属于另一类图。第~\ref{chap:rs-retrieval}~章会利用这些关系形成目录候选；本章先要求记录清楚边的含义和取得时间，避免不同关系被统一写成不明来源的“相似”。

一次请求也可能属于一个群组。共同选电影时，应知道哪些成员参加、哪个反馈属于共同选择，不能把某位成员事后的个人观看混入群组标签。个人与群组模型都需要在请求时点$t$截断信息。以下用$H_{u,t}=(e_1,\ldots,e_n)$表示已知历史，事件$e_k$至少包含物品$i_k$、动作$b_k$及时间$\tau_k$。浏览、收藏、加购与购买可以共同作为输入，但它们的标签和预测窗口须分别定义。

贯穿本章的商品例子只用于检查计算。我们会明确给出假设向量、权重和时间，不把算出的兴趣表示当作用户真实心理，也不把构造数据当作论文实验。固定相同事件再改变编码方式，才能看清方法差别来自顺序、候选选择还是额外信息。

% Outline: 2.2
\section{需求与对象建模}
\label{sec:rs-understanding-models}

需求、对象和匹配是本节的三个观察角度。用户模型回答哪些历史或表达应影响当前判断，对象模型回答有哪些可用事实及关系，匹配模型回答这些信息怎样共同支持选择。实际网络可以联合训练三者，不必划成三个独立服务；但讲解时仍要分清中间表示与最终预测，才能判断某个模块是否可以移到召回、排序或其他任务中。

% Outline: 2.2.1
\subsection{用户需求理解}
\label{sec:rs-user-understanding}

当前表达通常比一份长期画像更贴近这次任务。“只看500元以内”是明确限制，“轻便一些”需要结合场景解释，“不要这款”还需知道拒绝的是价格、款式还是已经购买。时间、地点、设备及刚刚进行的操作也会改变需求。历史可以填补未表达的部分，遇到冲突时则应保留表达的优先级或发起澄清。查询怎样对应搜索相关性，见第~\ref{chap:rs-search}~章。

行为建模面对的几组选择不能混为互斥类别。一个模型可以同时使用长历史、多兴趣表示和候选相关聚合。单一表示与多兴趣表示比较的是保留多少个兴趣方向，候选无关与候选相关比较的是候选在哪一步进入计算，短会话与长历史比较的是可用证据范围。例如，登山与摄影可以分别保留两个向量，再由多个向量取得候选；也可以保留完整历史，在评分相机时选择摄影记录。前一种目录接口在第~\ref{chap:rs-retrieval}~章的多兴趣模型中展开，后一种候选相关选择在第~\ref{chap:rs-ranking}~章基于历史行为的评分中展开。

训练目标同样不能由输入形式推断。使用购买和收藏历史，不表示两个事件具有相同价值；预测下一次浏览，也不同于预测给定曝光是否点击。表~\ref{tab:rs-behavior-models}先比较本章展开的候选无关表示模型的完整接口，候选相关模型则在第~\ref{chap:rs-ranking}~章基于历史行为的评分一节展开。辅助行为怎样支持多个预测任务，见第~\ref{chap:rs-ranking}~章；这些预测目标如何服务业务取舍，见第~\ref{chap:rs-constraints}~章。

这些模型共享两种常用的样本组织。下一物品任务将每个可见前缀与真实后继配对，从未交互目录采样对照；曝光响应任务则把请求前历史、已经展示的候选和随后观测到的响应组成样本。后者通常使用二元交叉熵
\begin{equation}
  \mathcal{L}_{\mathrm{click}}(\theta)
    =-\frac{1}{|\mathcal{D}|}\sum_{(H,i,y)\in\mathcal{D}}
       \left[y\log\widehat p_\theta(H,i)
         +(1-y)\log(1-\widehat p_\theta(H,i))\right],
  \label{eq:rs-understanding-click}
\end{equation}
其中$\mathcal{D}$为曝光样本集，$y\in\{0,1\}$为规定窗口内的点击标签。正项惩罚漏报点击，负项惩罚给未点击曝光过高的概率；对两项求梯度，才能让历史聚合朝实际预测目标调整。若另采样负曝光，训练样本中的正负比例已改变，排序效果与概率校准须分别检查。后文说明每个模型更新哪些模块、增加什么辅助目标，不把同一个点击损失重复当作架构创新。

在学习参数之前，也有直接利用历史统计的会话基线。一阶转移模型保存“刚点击物品到下一物品”的计数或概率；会话近邻先把当前会话与历史会话比较，再按邻近会话中出现的物品加权评分。前者只需保留末次物品，后者需要当前会话集合和可查询的历史会话索引。它们不必训练深层编码器，但新事件与新目录仍需更新计数或索引。下文神经模型的作用应通过这样的任务基线检验，而不是由网络层数判定；第~\ref{sec:rs-understanding-quality}~节给出同数据下方向不同的比较结果\scite{rs-ch02-session-evaluation}

\begin{table}[htbp]
  \centering
  \small
  \renewcommand{\arraystretch}{1.15}
  \begin{tabularx}{\linewidth}{@{}lXXXX@{}}
    \toprule
    模型 & 输入及组织 & 候选参与位置 & 中间输出 & 监督与应用输出 \\
    \midrule
    GRU4Rec & 匿名会话，逐事件更新 & 状态编码后 & 会话循环状态 & 下一物品比较；目录得分 \\
    SASRec & 物品前缀，因果注意 & 历史编码后 & 末端历史向量 & 后继正负例；目录分数 \\
    BERT4Rec & 已知历史，掩码重建 & 历史编码后 & MASK位置向量 & 遮蔽物品；目录概率 \\
    SR-GNN & 会话内转移图 & 会话编码后 & 局部与整体会话向量 & 下一点击；目录得分 \\
    MIMN & 增量行为、固定记忆 & 预测网络 & 记忆及归纳状态 & 点击与记忆正则；CTR \\
    PinnerFormer & 多时点未来窗口标签 & 用户编码后 & 批量用户向量 & 窗口内正反馈；目录得分 \\
    TiSASRec & 真实时间间隔 & 历史编码后 & 末端历史向量 & 后继正负例；目录分数 \\
    \bottomrule
  \end{tabularx}
  \caption{候选无关的行为表示模型。它们先把历史压成可复用的用户状态或向量，再由召回或排序与候选比较；给定候选后才激活、检索历史的模型在第~\ref{chap:rs-ranking}~章基于历史行为的评分中展开。}
  \label{tab:rs-behavior-models}
\end{table}

% Outline: 2.2.1.1
\subsubsection{GRU4Rec：循环状态的会话建模}
\label{sec:rs-gru4rec}

匿名访问没有可用的长期用户向量，而前序的一阶转移模型只保存“刚点击物品到下一物品”的计数，会话近邻又要把整段会话与历史会话逐一比较；二者都没有把会话内已经发生的多步顺序累积成一个紧凑依据。\term{GRU4Rec}的主要想法是用门控循环单元（gated recurrent unit, GRU）随事件逐步更新一个会话状态：新点击经过门控决定写入多少，再由当前状态为下一物品计算得分。原模型以物品标识的独热向量作为输入；论文也考察了历史独热向量的衰减聚合和附加嵌入层，不能将后来常见的“先查嵌入再输入GRU”说成其唯一原始设计\scite{rs-r10}

\textbf{模型结构。}第一步是输入：设$\bm{x}_k\in\{0,1\}^{|I|}$对应会话第$k$次点击，只在被点物品对应的位置为1，隐藏状态为$\bm{h}_k\in\mathbb{R}^d$。第二步是门控更新，为了看清新事件怎样改变依据，可把GRU更新写成
\begin{equation}
  \bm{h}_k=(\bm{1}-\bm{z}_k)\odot\bm{h}_{k-1}
       +\bm{z}_k\odot\widetilde{\bm{h}}_k,\qquad
  s_{k,i}=f_{\mathrm{out},i}(\bm{h}_k),
  \label{eq:rs-gru4rec-state}
\end{equation}
其中更新门$\bm{z}_k$和提议状态$\widetilde{\bm{h}}_k$由当前独热输入、此前状态及重置门共同计算，$f_{\mathrm{out}}$是输出层，把状态映射到目录物品的logit$s_{k,i}$。这里采用“门越大，新状态写入越多”的记号；另一种GRU约定可把门取补，含义不变。第三步是评分：对候选物品$i$，以输出层计算$s_{k,i}$，目录排名即比较各物品得分。候选物品只在这一步与状态比较，不进入状态递推本身，因此会话状态可以随行为增量维护。

训练要让实际下一物品$i^+$超过采样物品$j$。令$J_k$为本步负例集合，$\sigma$为sigmoid，原文讨论的两项损失为
\begin{align}
  \ell_{\mathrm{BPR}}
    &=-\frac{1}{|J_k|}\sum_{j\in J_k}\log\sigma(s_{k,i^+}-s_{k,j}),\\
  \ell_{\mathrm{TOP1}}
    &=\frac{1}{|J_k|}\sum_{j\in J_k}
       \bigl[\sigma(s_{k,j}-s_{k,i^+})+\sigma(s_{k,j}^2)\bigr].
  \label{eq:rs-gru4rec-top1}
\end{align}
两项损失的分工不同：BPR只要求正例得分高于每个采样负例，TOP1的第一项平滑近似负例排在正例前的程度，第二项再把负例得分向零约束。共同给正负分数加同一常数不会改变BPR第一项；原论文担心的是某物品在一条样本中作为正例被抬分，在其他样本中又充当负例，反复更新导致分数上漂，因此增加对负例绝对尺度的约束。未观测物品在这里是训练对照，不是已证实不喜欢的物品。

用一个教学例子检查状态更新。假设点击登山包后的状态为$(0.6,0.1)$，水壶事件产生提议状态$(0.8,0.2)$和门$(0.5,0.5)$，按式~\eqref{eq:rs-gru4rec-state}，新状态为$(0.7,0.15)$。若教学输出层对水壶配件与相机分别采用向量$(1,0)$和$(0,1)$，得分就是0.7和0.15。状态改变了下一次比较，却没有建立跨会话记忆：新会话重置后，这段个体信息不会自行保留。

训练沿会话展开反向传播，更新GRU的门控矩阵和输出层；采样输出只更新本批涉及的物品权重。会话并行批次在每个槽位放一段会话，依次处理各会话的下一事件，某个会话结束后换入新会话并重置该槽位的状态；原方法利用同批其他样本的目标物品作为负例，节省输出计算，但也使采样频率与物品流行度相关。隐藏维度、批次大小与dropout率由验证集选择，不是每个用户独立拟合的参数。服务端保存当前会话状态及已训练权重，新点击只执行一次递推；会话状态属于可随行为增量维护的候选无关表示，候选只在末端与状态比较，不需要对每个候选重新读取历史。没有当前会话时，不能把零状态预测解释为已有个体偏好。

原论文在电商RSC15和视频VIDEO上逐事件预测下一项。RSC15训练含约796万会话、3164万点击，测试用随后一天的会话，排名覆盖37,483件训练物品；VIDEO训练约300万会话、1,300万观看事件，但测试仅与最热门30,000个视频比较。表3中1,000维TOP1模型在RSC15的Recall@20为0.6206、MRR@20为0.2693，Item-KNN分别为0.5065、0.2048；VIDEO的BPR-1,000维模型虽提高召回，MRR却从Item-KNN的0.3381降为0.3136。增加深度、加额外前馈层和独立物品嵌入没有带来同样收益。这些是匿名电商和视频会话的离线证据，既不等同于曝光CTR，也不构成线上收入提升的测量\scite{rs-r10}

GRU4Rec把全部前缀压在一个递推变量中，预测时无法直接回看第几个位置发生了什么，这一缺口由下一节SASRec的自注意力直接读取各位置来补充；会话内出现回访和分支时，单条递推链也表达不出重复转移，SR-GNN为此把会话建成转移图。三者都属于先形成候选无关表示、再与候选比较的接口，与第~\ref{chap:rs-ranking}~章给定候选后才激活历史的评分模型不同。
% Outline: 2.2.1.2
\subsubsection{SASRec：自注意力的历史序列建模}
\label{sec:rs-sasrec}

GRU4Rec把全部前缀压在一个递推变量中，中间位置的影响被后续更新不断覆盖，预测时无法直接回看第几个位置发生了什么；若登山与摄影事件交错，需要让各位置直接参与当前表示。\term{SASRec}的主要想法是用自注意力替代递推：将物品嵌入与位置嵌入相加，以因果自注意力编码前缀，再用末端表示与物品向量比较。候选物品参与最后的评分，不进入历史编码本身，因此历史编码与候选无关\scite{rs-r11}

\textbf{模型结构。}第一步是输入：设截断后的历史长$n$，输入矩阵为$\bm{X}\in\mathbb{R}^{n\times d}$，其第$k$行为物品嵌入$\bm{v}_{i_k}$与位置嵌入$\bm{p}_k$相加。第二步是因果自注意力，一层的核心计算是
\begin{equation}
  \bm{A}=\operatorname{softmax}_{\mathrm{row}}
  \left(\frac{\bm{Q}\bm{K}^{\top}}{\sqrt d}+\bm{M}\right),
  \qquad \bm{Z}=\bm{A}\bm{V},
  \label{eq:rs-sasrec-attention}
\end{equation}
其中$\bm{Q},\bm{K},\bm{V}$由输入线性投影得到。掩码$\bm{M}$把未来位置及填充位置排除；填充不应因为恰好有零向量而获得一份注意权重。第三步是前馈与残差：位置前馈变换增加非线性，残差保留原表示，层归一化和dropout支持训练，堆叠多层后取末端行作为$\bm{h}_n$。通用运算见\BookChapterRef{chap:transformer}{模型卷的“Transformer”章}。

训练时每个合法位置$k$用$i_{k+1}$作正例，并从目录采样$j_k$作对照，最小化
\begin{equation}
  \ell_k=-\log\sigma(\langle\bm{h}_k,\bm{v}_{i_{k+1}}\rangle)
  -\log\bigl[1-\sigma(\langle\bm{h}_k,\bm{v}_{j_k}\rangle)\bigr].
  \label{eq:rs-sasrec-loss}
\end{equation}
推断取$\bm{h}_n$，为目录物品计算内积。训练位置$k$不能看到$i_{k+1}$，即使这个物品已存在于整条训练序列中；因果掩码正是为了让一次前向计算同时给出各位置的前缀预测，仍符合各自的预测条件。

假设历史为登山包、水壶、相机，最后一行注意权重为$(0.2,0.1,0.7)$，值向量依次为$(1,0),(0.8,0.2),(0,1)$，则该层聚合结果为$(0.28,0.72)$。这一计算说明末端可以保留多个先前位置的影响，不能据此宣称任意相机候选都会再次改变该组权重。长度为$n$的稠密注意力仍需计算约$n^2$个位置关系；截断会遗漏窗口外事件，扩窗则增加成本。完整目录排名和少量采样候选排名也不是同一难度，评价时须分别报告。

各位置的损失共同反传到共享物品表、位置表、注意投影和前馈网络，原实现用Adam训练，以验证指标连续20轮不改善作为停止条件。输入与评分共用物品表，避免另学一套同目录向量。由于历史编码不依赖候选，$\bm{h}_n$与物品向量的内积可以在编码完成后对整个目录计算，物品表可预先保存；但普通实现对每个新前缀重新编码全部位置，不像GRU那样单步递推更新，不能仅因使用因果掩码就假定已经提供增量缓存。

原论文在Amazon Beauty、Games、Steam和MovieLens-1M上保留用户末次事件测试、倒数第二次验证，并将测试正例与100个随机负物品排名。MovieLens-1M的历史窗口为200，其余为50；表III中Steam的Hit@10从GRU4Rec+的0.8018升至0.8729，NDCG@10从0.5595升至0.6306。但消融表IV显示，固定50维时Beauty移除位置嵌入的NDCG@10为0.3183，高于默认0.3142；加入双头注意也未超过单头。其实际研究任务包括商品、游戏与电影偏好预测，结果支持按稀疏程度选择组件，不能推出顺序越精细或头数越多越好\scite{rs-r11}

SASRec的每个位置只能读取左侧前缀，监督信号是紧邻的下一物品；已经完整观察到的一段历史还可以从右侧提供线索，BERT4Rec为此改用掩码重建的双向编码。两者的差别首先在监督方向与可见范围，而不是网络层数。
% Outline: 2.2.1.3
\subsubsection{BERT4Rec：掩码训练的双向历史建模}
\label{sec:rs-bert4rec}

SASRec逐位置预测下一物品时，每个位置只能读取左侧前缀，右侧已经发生的事件没有进入该位置的监督。已经完整观察到的一段历史还可以提供另一种练习：遮住其中某个物品，让模型根据其余已知事件恢复它。\term{BERT4Rec}的主要想法是采用这种掩码物品重建（masked item prediction），以双向Transformer编码物品与位置嵌入，让每个内部位置同时读取左右两侧的已知事件。它借鉴语言建模任务的形式，但编码的词表是物品目录，并非直接把NLP中预训练好的BERT当作推荐器\scite{rs-r26}

\textbf{模型结构。}第一步是构造输入：令$S$为被选作重建目标的位置集合，选中位置替换为MASK标记，$\widetilde H$为经过遮蔽的输入，物品嵌入与位置嵌入相加后送入双向Transformer。第二步是重建概率，$\bm{h}_k$为第$k$个输出，通过变换后的$\bm{h}_k$和物品输出嵌入得到目录概率，训练损失为
\begin{equation}
  \ell_{\mathrm{mask}}
    =-\frac{1}{|S|}\sum_{k\in S}
        \log p_\theta(i_k\mid\widetilde H).
  \label{eq:rs-bert4rec-mask}
\end{equation}
由于目标物品已经被隐藏，内部位置可以读取左右两侧的其他已知物品。第三步是推断：预测下一物品时，在请求前历史末端追加MASK，以该位置输出分布；训练还加入末项遮蔽情形，使模型接触这一使用接口。

对“登山包、MASK、相机”的重建，答案可能是曾被遮蔽的水壶；对“登山包、水壶、相机、MASK”的推断，目标则是尚未发生的下一行为。假设后一种输入给镜头、水壶、背包配件的概率为0.5、0.2、0.3，则本次排序为镜头、配件、水壶。前一个训练问题的答案不能直接搬到后一个问题上。

双向的范围是请求前已知历史。训练切分、掩码位置及末项构造一旦把请求后的行为混入，就改变了可用证据。完整目录softmax还可能昂贵；减少训练候选需要重新说明目标近似和评价口径。模型擅长何种历史条件，应由同一信息边界下的比较确定。

每轮重新采样遮蔽位置，使同一历史产生不同恢复练习。只有遮蔽位置进入重建损失，梯度却会经过它读取的全部已知上下文，更新物品与位置嵌入、双向注意层和输出变换。原论文直接将选中物品替换为MASK，不应自行补入语言BERT的其他替换比例；其Adam学习率为$10^{-4}$，遮蔽率按数据集验证选择。推断保留物品词表、模型参数和已知历史，取消训练随机遮蔽，只追加末端MASK；双向编码在请求前完成，末端MASK位置的输出再与物品向量比较，候选不进入编码，这一接口与SASRec相同，差别只在训练时的可见方向与监督目标。

原实验覆盖Beauty、Steam、MovieLens-1M和20M，测试候选为一个真实后继加100个按流行度采样的未交互物品。表2中MovieLens-1M的HR@10为0.6970，对照SASRec为0.6629；表3将每次遮蔽限制为一个位置后得到0.6869，说明原设置下双向读取和更多重建目标各有贡献。电影窗口200、Beauty和Steam窗口50，不能把此表与SASRec原论文的不同预处理结果直接相减。后续同代码、全目录比较又显示训练损失改变足以逆转多数数据集的模型次序，详见第~\ref{sec:rs-understanding-quality}~节；因此双向重建是一种学习选择，不是下一物品预测的无条件优势\scite{rs-r26,rs-ch02-gold-loss}

SASRec与BERT4Rec都把会话看作一条带位置的序列；当同一物品在会话中重复出现、转移形成回访和分支时，序列位置表达不出两个后继各占多少比例，SR-GNN为此把一次会话建成物品转移图。
% Outline: 2.2.1.4
\subsubsection{SR-GNN：物品转移图的会话建模}
\label{sec:rs-srgnn}

GRU4Rec与SASRec都把会话当作一条线性位置序列；当同一物品在会话中重复出现、从它出发又走向两个不同后继时，序列位置既合并了回访，也表达不出两条后继各占多少比例。\term{SR-GNN}的主要想法是把一次会话的不同物品建成节点，以相邻点击构造有向转移边，用图结构保留回访与分支，再通过门控图传播整合入边与出边信息，组合末次物品和整个会话的表示来预测下一次点击\scite{rs-r13}

\textbf{模型结构。}第一步是建图：会话中出现过的不同物品各为一个节点，相邻两次点击构成一条有向转移边，重复转移按固定规则计数或合并。第二步是门控图传播：设节点$i$的表示为$\bm{v}_i^{(\ell)}$，$\bm{A}^{\mathrm{in}}$和$\bm{A}^{\mathrm{out}}$按相应连接关系归一化；一次传播先分别汇入前驱和后继表示，将两侧线性变换的结果拼接成$\bm{a}_i^{(\ell)}$，再以GRU式门更新节点。重复转移的计数规则须在出入邻域内归一化，且训练和服务保持一致。图~\ref{fig:rs-srgnn-graph}给出一个会话的建图结果。第三步是组合局部与全局：末次点击$i_n$给出局部表示$\bm{s}_{\mathrm{local}}=\bm{v}_{i_n}^{(L)}$，会话整体表示为
\begin{equation}
  \bm{s}_{\mathrm{global}}=\sum_{k=1}^{n}\alpha_k\bm{v}_{i_k}^{(L)},
  \qquad
  \bm{s}=\bm{W}
  [\bm{s}_{\mathrm{local}};\bm{s}_{\mathrm{global}}],
  \label{eq:rs-srgnn-session}
\end{equation}
其中$\alpha_k=\bm{q}^{\top}\sigma(\bm{W}_1\bm{v}_{i_n}^{(L)}
+\bm{W}_2\bm{v}_{i_k}^{(L)}+\bm{c})$由最后事件和各事件共同计算。原形式的$\alpha_k$不是必须归一化到和为一的概率。最后用$\langle\bm{s},\bm{v}_i\rangle$形成物品logit，对实际下一点击的softmax交叉熵进行训练。

\begin{figure}[htbp]
  \centering
  % SR-GNN：单个匿名会话的物品转移图（教学示例：登山包、水壶、登山包、相机）。
% 单栏图（figure），内容宽约78mm；回访合并为同一节点，出边按频次归一化。
\begin{tikzpicture}[x=1mm,y=1mm,
  font=\figurefont\footnotesize,
  gnode/.style={rectangle,draw=FigureModel!70,fill=FigureModel!10,align=center,
    text=BookInk,rounded corners=1.4pt,minimum height=9mm,
    text width=16mm,inner sep=2pt},
  glocal/.style={rectangle,draw=FigureOutput!70,thick,fill=FigureOutput!14,align=center,
    text=BookInk,rounded corners=1.4pt,minimum height=9mm,
    text width=16mm,inner sep=2pt},
  gflow/.style={-{Stealth[length=1.8mm,width=1.2mm]},draw=FigureOutput!85,line width=.85pt},
  gback/.style={-{Stealth[length=1.8mm,width=1.2mm]},draw=FigureModel!80,line width=.8pt,dashed},
  glab/.style={font=\figurefont\scriptsize,text=BookMuted,fill=white,inner sep=1pt}
]
  % 节点：回访的登山包只建一次
  \node[gnode] (a) at (16,20) {登山包\\（回访合并）};
  \node[gnode] (b) at (58,33) {水壶};
  \node[glocal] (c) at (58,9) {相机\\（末次局部）};

  % 转移边：包→水壶，水壶→包（回访），包→相机
  \draw[gflow] (a.north east) .. controls (36,34) .. (b.west);
  \node[glab,align=center] at (37,30) {出边 $1/2$};
  \draw[gback] (b.south west) .. controls (36,25) .. (a.north east);
  \node[glab,align=center] at (36,22) {回访};
  \draw[gflow] (a.south east) .. controls (38,12) .. (c.west);
  \node[glab,align=center] at (38,9) {出边 $1/2$};

  % 局部与全局说明
  \node[glab,align=left] at (16,-2)
    {局部：末次节点 $c$\\全局：$\alpha_k$ 对全部节点聚合};
\end{tikzpicture}

  \caption{SR-GNN的会话转移图（教学示例）。序列“登山包、水壶、登山包、相机”中，登山包只建一个节点并拥有两条等权出边，回访表现为节点重复出现而非新增节点；末次相机作为局部兴趣，全局注意力在全部节点上聚合。图只表达本会话的转移，不含其他用户的协同关系。}
  \label{fig:rs-srgnn-graph}
\end{figure}

会话“登山包、水壶、登山包、相机”产生三条转移：包到水壶、水壶到包、包到相机。包的两个不同后继在等权情况下各占出边归一化权重的$1/2$；末次节点是相机，整体聚合中登山包又出现两次。若再发生一次包到水壶，采用转移频次的图会改变这两个后继的比例。图因此能表达回访结构，但只保留节点、边和用于聚合的事件信息，不能从邻接矩阵唯一还原完整事件顺序。它也没有包含其他用户的协同关系。

训练把一次会话拆成多个前缀—后继样本，对每个前缀重新建图；反向传播同时更新物品嵌入、入出边变换、门控传播、会话注意与输出投影。原作者实现使用目录logit的交叉熵，论文印出的式10另含非目标项，复现时应明确采用哪个目标，不能无声将二者视为同一公式。会话图在当前会话内维护，候选只在末端与$\bm{s}$做内积比较，不进入图传播；服务端保留当前会话及模型参数，新事件可能改变节点和转移边，需要相应更新图传播，而不是只追加一个孤立节点。图只覆盖本次会话，新会话重新建图，不保留跨会话记忆。

Yoochoose与Diginetica实验删除单事件会话及低频物品，用随后一天或一周作测试；Yoochoose另外比较最近的$1/64$和$1/4$训练切片。表2中Diginetica的前20项命中比例为50.73\%、MRR@20为17.59\%，NARM为49.70\%、16.17\%；原文将前者命名为P@20，但在单后继任务中它是命中比例，不能理解为每个列表20项中有一半相关。图结构消融中，引入全局会话边或会话内全连接未持续超过原转移图。这些电商日志证据支持局部转移建模，未提供生产A/B业务收益\scite{rs-r13}

以上GRU4Rec、SASRec、BERT4Rec与SR-GNN都只覆盖当前会话或截断窗口内的事件；每次请求重新扫描完整历史，成本随事件数持续增加，MIMN为此把压缩工作移到行为到来时，用固定容量记忆延长可用历史。
% Outline: 2.2.1.8
\subsubsection{MIMN：增量记忆的长历史压缩}
\label{sec:rs-mimn}

前面的序列模型要么只覆盖当前会话，要么把历史截断到几百条；每次请求都重新扫描完整历史，成本会随事件数持续增加，而候选需要的久远细节可能恰好在截断之外。\term{MIMN}（multi-channel interest memory network）的主要想法是把一部分工作移到行为到来时：维护固定数量的记忆槽，新事件读取和修改记忆，预测请求再使用已经维护的状态。网络由基础读写记忆和记忆归纳单元共同组成\scite{rs-r27}

\textbf{模型结构。}第一步是控制器寻址：设记忆矩阵$\bm{M}^{(k)}\in\mathbb{R}^{m\times d}$包含$m$个槽，控制器由新行为生成读取键、写入键、擦除向量$\bm{e}_k$和添加向量$\bm{a}_k$；内容寻址将键与各槽的余弦相似度归一化为读取或写入权重。第二步是读出与擦除写入，用$\bm{w}^{\mathrm{r}}_k,\bm{w}^{\mathrm{w}}_k\in\mathbb{R}^m$表示这些权重，基本操作为
\begin{align}
  \bm{r}_k&=\sum_{j=1}^{m}w^{\mathrm{r}}_{k,j}\bm{M}^{(k-1)}_{j,:},\\
  \bm{M}^{(k)}_{j,:}
    &=\bm{M}^{(k-1)}_{j,:}\odot
        (\bm{1}-w^{\mathrm{w}}_{k,j}\bm{e}_k)
       +w^{\mathrm{w}}_{k,j}\bm{a}_k.
  \label{eq:rs-mimn-write}
\end{align}
写入因此既可减弱旧内容，也可加入新信息。模型还根据累计使用量重新平衡写入权重，并惩罚各槽累计写入量的方差，以减少少数槽长期承受几乎全部更新的情况。图~\ref{fig:rs-mimn-memory}给出一次写入的数据路径。第三步是记忆归纳：单纯槽值还不足以表示更新轨迹，归纳单元为槽位维护多通道GRU状态$\bm{S}$，选择读取权重较高的通道，结合相应记忆与当前行为更新，各通道共享GRU参数。读出及归纳状态与候选、用户和情境一起进入CTR网络，训练由点击损失和记忆使用正则共同约束。

\begin{figure}[htbp]
  \centering
  % MIMN：新行为经控制器寻址固定记忆槽，读出并擦除写入；候选只在CTR读出阶段进入。
% 单栏图（figure），内容宽约80mm；上行为新行为→控制器，记忆矩阵在下方，读出标注独立置于框下。
\begin{tikzpicture}[x=1mm,y=1mm,
  font=\figurefont\footnotesize,
  mbeh/.style={rectangle,draw=FigureInput!70,fill=FigureInput!12,align=center,
    text=BookInk,rounded corners=1.4pt,minimum height=9mm,text width=18mm,inner sep=2pt},
  mctl/.style={rectangle,draw=FigureModel!70,thick,fill=FigureModel!10,align=center,
    text=BookInk,rounded corners=1.4pt,minimum height=16mm,text width=24mm,inner sep=2pt},
  mmem/.style={rectangle,draw=FigureLoss!75,fill=FigureLoss!12,align=center,
    text=BookInk,rounded corners=1.4pt,minimum height=26mm,text width=18mm,inner sep=2pt},
  mout/.style={rectangle,draw=FigureOutput!70,thick,fill=FigureOutput!14,align=center,
    text=BookInk,rounded corners=1.4pt,minimum height=11mm,text width=20mm,inner sep=2pt},
  mflow/.style={-{Stealth[length=1.8mm,width=1.2mm]},draw=FigureOutput!85,line width=.85pt},
  mwrite/.style={-{Stealth[length=1.8mm,width=1.2mm]},draw=FigureLoss!85,line width=.85pt},
  mlab/.style={font=\figurefont\scriptsize,text=BookMuted,fill=white,inner sep=1pt}
]
  % 上行：新行为 → 控制器
  \node[mbeh] (beh) at (12,26) {新行为 $e_k$};
  \node[mctl] (ctl) at (42,26)
    {控制器\\读键 $w^{\mathrm r}_k$／写键 $w^{\mathrm w}_k$\\擦除 $\bm e_k$／添加 $\bm a_k$};

  % 记忆矩阵（与控制器留出间隙，仅含槽位说明，无叠加文字）
  \node[mmem] (mem) at (42,0) {记忆矩阵 $\bm M$\\$m$ 个槽位};
  % CTR 与记忆矩阵同基线，箭头从右侧水平引出
  \node[mout] (ctr) at (70,0) {CTR 网络\\候选在此进入};

  \draw[mflow] (beh.east) -- (ctl.west);
  \draw[mwrite] (ctl.south) -- (mem.north);
  \node[mlab,align=left] at (49,15.5) {擦除＋写入};
  \draw[mflow] (mem.east) -- (ctr.west);

  % 读出标注独立置于记忆矩阵下方，不压框内文字、不压箭头
  \node[mlab,align=center] at (42,-18) {读出 $\bm r_k$ 与归纳状态，进入 CTR};
\end{tikzpicture}

  \caption{MIMN的记忆读写（教学示意）。新行为经控制器生成读键、写键、擦除与添加向量；内容寻址得到槽权重，先读出$\bm r_k$，再按擦除—写入更新记忆矩阵。记忆在行为到来时维护，候选只在CTR读出阶段进入，写入时并不知道未来评估哪个候选。}
  \label{fig:rs-mimn-memory}
\end{figure}

假设两个槽原为$(1,0)$与$(0,1)$，新摄影收藏的写入权重为$(0.2,0.8)$，擦除向量为$(0.5,0.5)$，添加向量为$(0,1)$。按式~\eqref{eq:rs-mimn-write}，新槽为$(0.9,0.2)$与$(0,1.4)$。同一事件改变了两个槽，幅度不同；槽的含义来自学习，不能预先断言每个槽始终对应一个固定兴趣类别。

原方案中的用户兴趣中心把这些更新放在行为触发路径，服务请求读取状态，减少重复历史计算。记忆槽的维护不依赖候选，可以在行为到来时增量完成并预存；但候选在CTR读出阶段才进入，逐候选的响应计算仍在请求时完成，记忆本身不能替代这一步。固定$m$意味着在线读写规模不直接随已发生事件总数增长，但状态存储、行为更新和版本同步仍有代价，历史细节也可能丢失。新行为是否已经写入，属于第~\ref{sec:rs-evidence-refresh}~节的时效问题；能够持续写入多年历史，不等于无损保存多年事件。

训练仍沿历史展开记忆更新并反向传播，更新控制器、寻址、共享归纳GRU和预测网络；上线把这条展开路径改为事件驱动的增量执行，而不是上线后自行优化每个用户的记忆参数。原部署先离线回放120天行为，把积累后的槽值和归纳状态推送到兴趣中心，避免从空记忆开始长时间预热。参数更新与用户状态写入各有版本，回滚还需恢复未受污染的状态；仅更换CTR头不能撤销错误事件已经写入的记忆。

Amazon Books实验过滤短于20的历史并截断100条，Taobao截断200条，工业数据使用60天历史并截断1,000条。表4中Taobao基础记忆AUC为0.9070，加使用量正则为0.9112，加归纳单元为0.9179；比较所用DIEN去掉了辅助损失，不能当作完整DIEN的绝对胜负。槽数消融在Books以4槽最好、Taobao以8槽最好。工业表5中MIMN为0.6644、DIEN为0.6541，加入大促日训练数据反而降至0.6627。2019年3月30日至5月10日广告A/B相对原线上DIEN报告CTR提高7.5\%、RPM提高6\%，并给出兴趣中心部署经验；其中增加可用历史与改变请求计算同时发生，收益不是固定输入下仅换记忆模块的净效果\scite{rs-r27}

MIMN在写入记忆时并不知道未来会评估哪个候选，固定容量的压缩会把候选所需的细节混掉；保留原始长历史、先按候选检索再精细匹配的思路，在第~\ref{chap:rs-ranking}~章的SIM中展开。PinnerFormer处理的是另一个方向的问题：批量用户向量的监督范围不应只限于紧邻的下一动作。
% Outline: 2.2.1.15
\subsubsection{PinnerFormer：多时点目标的长期兴趣表示}
\label{sec:rs-pinnerformer}

下一物品目标只监督紧邻动作，而一份批量更新、供多次请求使用的用户向量，还要概括未来一段时间的兴趣，只按“紧接着会点什么”训练并不匹配这个用途。\term{PinnerFormer}的主要想法是把训练监督扩展到未来一段时间中的多个正反馈，让同一表示概括窗口内的兴趣。它延长的是预测范围，尚未估计某次推荐行动对未来的影响\scite{rs-r78}

\textbf{模型结构。}第一步是输入与编码：输入包括物品内容向量、动作类型、入口、时间与持续时长等特征，因果Transformer编码各个已知前缀，通过投影及$\ell_2$归一化得到用户向量；物品侧也将内容向量映射到同一空间。第二步是多时点窗口对齐：训练在多个历史位置$k$选择未来窗口$P_k$中的正物品$i^+$，用采样softmax将$\bm{h}_k$与这些物品对齐。忽略采样校正时，一对正样本的目标为
\begin{equation}
  \ell_{k,i^+}=-\log
  \frac{\exp(\langle\bm{h}_k,\bm{v}_{i^+}\rangle/\tau)}
  {\exp(\langle\bm{h}_k,\bm{v}_{i^+}\rangle/\tau)
    +\sum_{j\in J_k}\exp(\langle\bm{h}_k,\bm{v}_j\rangle/\tau)}.
  \label{eq:rs-pinnerformer-loss}
\end{equation}
温度$\tau>0$控制分数差的尺度。原方案结合批内与随机负例，处理正物品误作负例，并对采样概率作logit校正；多个时点和窗口正例共同提供监督，而非让一个末端状态独自承担所有未来目标。

假设一次户外浏览之后，未来七天内用户收藏水壶、帐篷和摄影背包。下一动作目标只监督紧邻事件，窗口目标则可让三个收藏都约束同一前缀的表示。三者在训练时是未来标签，生成该前缀表示时都不可见。服务时批量编码用户，将向量用于近邻候选或下游评分特征。

窗口越长，标签越可能包含多个任务或变化的需求；是否合适取决于向量的使用周期和目标行为。向量一天更新一次也可能错过刚发生的兴趣变化，因而可以与第~\ref{sec:rs-transact}~节的即时分支组合。预测未来窗口的参与行为与控制未来参与行为之间，仍隔着行动选择和因果证据。

密集监督在多个前缀位置采样未来28天的正反馈，分别让梯度经过相应位置输出，避免把全部未来目标的梯度先平均到一个末端向量。批内负例偏向热门正反馈，随机目录负例分布不同，所以原方法使用混合池并作采样概率校正；与同用户窗口正例重合的负例需要排除。Adam更新历史Transformer及用户、物品投影，而作为输入的内容向量有自己的训练来源。用户向量与候选无关，可在离线批量阶段编码并预存，上线每天为有活动用户重新编码，保存一份归一化用户向量及目录向量，供近邻候选或下游评分特征使用；服务时不运行训练中的未来窗口采样。

Pinterest离线以未来14天的保存及长点击等正反馈评价，另以100万件物品索引测覆盖。P90覆盖表示占全部用户前10候选出现量90\%的物品集合占索引的比例，数值较小表示机会更集中。表3将SASRec式后继目标改成密集28天目标，Recall@10从0.198升至0.229，但P90目录覆盖从0.048降至0.042；表1的20兴趣PinnerSage覆盖更高，召回却更低，单向量并非在所有维度胜出。表4混合负例加校正提高召回，同时降低全局覆盖，说明负例配方与表示效果不能分开看。线上作为Homefeed排序特征替换PinnerSage，报告保存量提高7.5\%、DAU提高0.4\%；作为广告特征新增，则三个入口CTR提高7.1\%至10.0\%。这些是真实部署和A/B证据，但新增特征与替换特征的实验不同，也未证明28天监督优化了行动的长期因果价值\scite{rs-r78}

PinnerFormer改变的是监督覆盖的时间范围；TiSASRec则在同一后继目标下补充物品之间的真实时间间隔，让顺序相同但节奏不同的历史得到不同表示。
% Outline: 2.2.1.18
\subsubsection{TiSASRec：时间间隔条件下的序列建模}
\label{sec:rs-tisasrec}

SASRec的位置编号只给出事件次序，“登山包、水壶、相机”发生在半小时内还是三个月内，可能对应不同的判断，位置编码区分不了这种节奏。\term{TiSASRec}的主要想法是在SASRec式历史编码中增加两两真实时间间隔，并在注意的键和值两侧分别加入位置与间隔表示\scite{rs-r101}

\textbf{模型结构。}第一步是构造间隔索引：设请求前可见的最小非零时间间隔为$\delta_u$，截断上界为$r_{\max}$，位置$k,l$之间的间隔索引为
\begin{equation}
  r_{k,l}=\min\left(r_{\max},
     \left\lfloor\frac{|\tau_k-\tau_l|}{\delta_u}\right\rfloor\right).
  \label{eq:rs-tisasrec-time}
\end{equation}
全部时间戳相同时需为尺度规定兜底值；训练回放不能从请求后的事件估计$\delta_u$。第二步是条件化注意力：查询位置$k$对历史位置$l$的注意logit使用物品键、位置键及$r_{k,l}$的间隔键，聚合值也增加独立的位置值与间隔值；因果和填充掩码继续限制可见位置。第三步是评分：推断取末端表示与物品向量内积，候选不进入间隔编码。

假设时间戳为0、2、8小时，最小非零间隔为2小时，相对首项的缩放时间为0、1、4。若$r_{\max}=3$，两两查表索引为
\[
  \bm{R}=
  \begin{pmatrix}0&1&3\\1&0&3\\3&3&0\end{pmatrix}.
\]
末端对第二项读取间隔3，对自身读取0。若把末项改成4小时，最后一行改为$(2,1,0)$；物品顺序没变，但键和值中的时间信息都变了。

模型用各前缀的下一物品和采样负例作二元交叉熵训练，推断以末端表示和物品向量内积评分。为说明改变间隔怎样传到输出，假设物品值为$(1,0),(0.8,0.2),(0,1)$，两种时间输入产生权重$(0.2,0.2,0.6)$与$(0.2,0.5,0.3)$，暂把值侧附加项设零，则表示从$(0.36,0.64)$变为$(0.60,0.40)$，两个坐标方向候选的次序也随之反转。这只是给定参数结果的教学计算，时间变化并不必然造成该方向的变化。

按个人最小间隔缩放保留了相对节奏，也可能让绝对时长相差很大的两段历史拥有相同索引。截断则使大间隔共享表示。两项处理都需要按数据条件检验，时间权重本身不能证明过去行为对现在选择具有因果作用。

正负损失和正则共同更新物品表、位置键值表、间隔键值表及注意前馈网络，原配置两层、每前缀一个未交互负物品，采用Adam。最小间隔与截断索引来自数据处理而非梯度优化。推断保留物品与时间戳前缀并重新构造间隔矩阵，仍需保存位置表和两套间隔表；只缓存一个末端向量不能在新事件到来后恢复全部两两间隔。

六数据集实验用末项测试、倒数第二项验证，并与100个随机负物品排名。表5中MovieLens-1M同窗口SASRec的NDCG@10为0.5524，只有相对间隔的TiSASRec-R为0.5648，兼用位置与间隔为0.5706；但Steam只用间隔降至0.5731，低于SASRec的0.5801，二者结合才到0.5897。个人缩放也并非每个数据集最佳：Beauty不缩放为0.2825、个人缩放为0.2818。应用证据是电影、商品与游戏日志上的离线后继预测，论文没有用这些结果测线上点击或出行效果；时间信息是否足够精确，应先检查时间戳重复和粒度\scite{rs-r101}

以上模型的请求主体都是单个用户；当一次选择由一群人共同作出时，个人偏好如何聚合成群组依据成为新问题，AGREE与GroupIM依次处理成熟群组和临时稀疏群组。
% Outline: 2.2.1.20
\subsubsection{AGREE：群组请求下的成员偏好聚合}
\label{sec:rs-agree}

共同选一部电影时，个人推荐列表的交集可能为空，把成员偏好固定平均又未必能解释群组的实际选择：同一部电影对不同成员的重要性随候选变化。\term{AGREE}的主要想法是让成员权重随候选改变，结合成员聚合与群组自身嵌入预测共同交互，同时用个人交互训练共享表示\scite{rs-r103}

\textbf{模型结构。}第一步是成员注意：设群组$g$的成员集合为$U_g$，成员向量为$\bm{p}_u$、群组ID向量为$\bm{q}_g$，候选相关注意网络给出
\begin{equation}
  \alpha_{u,i}=
   \frac{\exp a_\theta(\bm{p}_u,\bm{v}_i)}
        {\sum_{v\in U_g}\exp a_\theta(\bm{p}_v,\bm{v}_i)},\qquad
  \bm{h}_g(i)=\sum_{u\in U_g}\alpha_{u,i}\bm{p}_u+\bm{q}_g.
  \label{eq:rs-agree}
\end{equation}
第二步是评分：预测网络读取$\bm{h}_g(i)$、$\bm{v}_i$及二者逐元素乘积，产生群组得分；个人路径用$\bm{p}_u$替换群组表示。两条路径共享成员和物品表示，以个人和群组的正负交互共同训练。原工作使用回归式成对目标，要求正负得分差接近1，可写成$(s_{g,i^+}-s_{g,j}-1)^2$，个人样本同理。

假设两名成员向量为$(1,0)$、$(0,1)$，某候选的权重为0.75、0.25，成员聚合为$(0.75,0.25)$；再加群组嵌入$(0.1,0.2)$，得到$(0.85,0.45)$。群组ID项提供了成员聚合之外的参数贡献。候选改变后权重也可能改变，服务时必须重新计算，不能把这一结果直接当作独立群组塔用于全部候选。

注意权重由预测损失学习，不是成员真实的表决权。未见过的新群组也没有已训练的$\bm{q}_g$，需要明确初始化或去掉该项后的适配方案。共同选择是否照顾每个成员，则属于第~\ref{chap:rs-constraints}~章的目标与约束，不能从群组预测准确率直接得到。

训练先以均匀成员权重预训练，再启用注意。每批混合个人和群组正交互，原配置每条正例采4个未交互物品，以成对目标更新成员表、物品表、群组ID表、注意网络和共享交互层。个人样本不更新不涉及的群组ID，却可改善群组读取的成员表示；群组样本则反向约束成员和候选权重。推断不再采负例，只输入实际成员与待比较物品，并取消训练dropout。若某用户没有个人记录但参加过训练群组，他仍可能由群组反馈取得表示，这与训练中完全未出现的新用户不同。

原研究在Mafengwo旅游群组和CAMRa2011家庭电影记录上比较，前者995个群组、3,595次共同交互，后者290个家庭、145,068次共同交互。每个用户或群组随机留一项测试，再与100个未交互物品排名，各配置重复5次。表2的Mafengwo群组HR@10为0.6400、NDCG@10为0.4244，均匀成员GREE为0.6151、0.4111；表4删除成员聚合或群组ID分别降至0.5709、0.6181的HR@10，说明两项并非重复。其电影评分全部转为正交互，不能把任务说成预测满意评分。后续GroupIM全目录、严格新群组实验给出另一边界：Weeplaces的AGREE Recall@20为0.484，VAE-CF均值聚合为0.502。这个结果不撤销AGREE原任务的观察，而是说明成熟群组ID与极稀疏新组合的条件不同；两篇都没有据这些离线表测真实群体满意度\scite{rs-r103,rs-r104}

AGREE的群组ID向量$\bm{q}_g$只对训练中反复出现的成熟群组有足够反馈；临时群组很少重复，为每个组合单独学一个ID向量数据不足，GroupIM因此放弃群组ID，改从成员各自的交互表示聚合。
% Outline: 2.2.1.21
\subsubsection{GroupIM：稀疏群组反馈下的偏好学习}
\label{sec:rs-groupim}

AGREE的群组ID依赖反复出现的成熟群组；临时群组很少重复出现，为每个组合单独学习一个ID向量可能没有足够数据。\term{GroupIM}的主要想法是从成员各自的交互向量编码个人表示，再通过均值、最大池化或注意聚合构造群组表示，利用成员关系辅助学习，而不要求每个新群组先拥有一份成熟ID嵌入\scite{rs-r104}

\textbf{模型结构。}第一步是成员编码与聚合：令$\bm{x}_u\in\{0,1\}^{|I|}$为成员交互向量，$\bm{p}_u=f_{\mathrm{enc}}(\bm{x}_u)$，$\bm{h}_g=f_{\mathrm{agg}}(\{\bm{p}_u:u\in U_g\})$。目录预测头输出$p_\theta(i\mid g)=\operatorname{softmax}(\bm{W}\bm{h}_g)_i$，以共同交互的多项似然训练。第二步是成员信息约束：成员历史不能一律等权当作共同偏好，所以模型另设双线性判别器
\begin{equation}
  D(u,g)=\sigma(\bm{p}_u^\top\bm{W}_{D}\bm{h}_g).
  \label{eq:rs-groupim-discriminator}
\end{equation}
真实成员与采样非成员构成对比样本，二元交叉熵推动用户—群组表示保留成员信息；这是互信息估计的训练约束，不是直接观测到群组协商过程。

判别分数还有第二个用途：以$D(u,g)$加权成员历史对群组目录预测的损失贡献。把常数归一化因子省略，该项为
\[
  \ell_{\mathrm{member},g}
    =-\sum_{u\in U_g}D(u,g)\sum_{i\in I}x_{u,i}\log p_\theta(i\mid g).
\]
训练交替进行：固定判别器，用共同反馈及加权成员反馈更新推荐路径；再训练对比目标，同时更新判别器、编码器和聚合器。个人交互可预训练编码器，服务阶段则直接编码成员历史、聚合并计算目录得分，不必运行负成员判别。

假设教学编码器给两名固定成员向量$(1,0)$与$(0,1)$。第三人是$(1,0)$时，均值群组向量为$(2/3,1/3)$；第三人换成$(0,1)$时，变为$(1/3,2/3)$。取$\bm{W}_D$为单位矩阵，第一名成员的判别值随之从$\sigma(2/3)\approx0.661$变为$\sigma(1/3)\approx0.583$，其历史监督贡献也改变。若目录头为单位映射，两个物品方向的概率从约$(0.583,0.417)$交换为$(0.417,0.583)$。

该例连起了成员变化、表示、判别权重和目录分布，但不能证明某人实际说服了其他成员。成员没有有效历史，或者群组内部存在强烈冲突时，辅助约束也无法凭空确定共同选择。

GroupIM的推荐损失是目录多项似然，不把每个未交互物品另当作一个人工负群组标签；负采样用于成员—群组判别，原配置每个真实成员配5个非成员。个人编码器先以个人交互预训练，然后参与推荐和互信息交替更新。上线只需编码请求前成员历史、按已训练聚合器构造群组向量，再访问目录，不需要保留判别器或为每个临时组合存ID向量。

原实验覆盖Yelp、Weeplaces、Gowalla和Douban，群组平均共同交互均少于3.5项。前三者以朋友在同地点15分钟内签到构造共同事件，Douban以活动及场地记录构造；群组分别随机分成70\%训练、10\%验证和20\%测试，成员个人历史都用于训练，测试是新成员组合而非全部新用户。表2中Weeplaces注意GroupIM的Recall@20为0.716，MoSAN为0.548；均值GroupIM为0.706，说明复杂注意并非核心收益的唯一来源。表3同数据仅群组损失的NDCG@50为0.420，加互信息为0.431，完整自适应成员监督为0.543。证据是实际签到和活动记录上的离线群组预测，群组构造仍可能误认共同决策，不能视为已经测量成员协商或满意度\scite{rs-r104}

与AGREE相比，GroupIM的群组向量由成员表示聚合得到，不依赖群组ID，因而能服务未见过的临时组合；代价是共同反馈稀疏时，聚合与判别都只能利用成员历史，无法引入群组长期积累的偏好。两者都只解决“把成员依据合成群组表示”，不测量真实协商过程。

候选相关的历史模型按各自任务放在使用它们的章节：DIN、DIEN、DSIN等短序列激活模型，以及SIM、TWIN、ETA、SDIM、TWIN-V2、LONGER等长历史选集与压缩模型，还有TransAct、HSTU和时空条件下的STAN，都在给定候选后才计算用户侧表示，归入第~\ref{chap:rs-ranking}~章基于历史行为的评分；多兴趣向量（MIND、ComiRec）在编码后用于访问目录，见第~\ref{chap:rs-retrieval}~章。表~\ref{tab:rs-history-tradeoffs}固定“历史已取得、候选需要评分”的任务，在第~\ref{chap:rs-ranking}~章比较这些方案请求时访问什么、行为到来时维护什么。

% Outline: 2.2.2
\subsection{对象信息理解}
\label{sec:rs-item-understanding}

用户表示再丰富，也无法替代对象事实。推荐一个通勤背包至少需要知道价格、容量、重量和防水说明；只有标题里的“户外”一词，既不能确认容量，也不能确认防水等级。对象理解把文本、图片、视频和元数据变成可使用的信息，同时保留哪些是直接登记、哪些由模型推断、哪些尚缺依据。

不同模态提供的证据范围不同。图片可以展示颜色、外观和可见结构，却通常不足以确认准确重量；结构化参数可能明确重量，却需要检查版本、计量单位及是否包含配件。语言和视觉编码器可以提取特征，预测价值仍须由当前任务检验。某个向量能识别背包主题，不表示它保留了区分通勤和重装登山所需的全部细节。

对象之间的关系也要标明含义。两个背包外形相似，不一定互相替代；相机与适配存储卡可能互补，但买相机和买水壶经常共现，也可能只是同一批旅游用户的消费习惯。品牌和主题关系来自登记知识，共现关系来自交互，视觉相似来自内容计算。关系能否承担兼容性、替代性或组合要求，取决于证据是否支持这种用途。

因此，对象资料不应只有一个无来源的“特征值”。对于重量0.8千克，可以保存厂家规格和更新时间；对于“适合通勤”，可以保存判定依据及不确定性；遇到详情页和结构化字段冲突，应保留冲突状态而不是由模型分数覆盖。下节将内容与关系真正接到匹配目标；第~\ref{chap:rs-presentation}~章还会检查展示表达有没有超出这些事实。

% Outline: 2.2.3
\subsection{需求与对象匹配}
\label{sec:rs-user-item-match}

匹配包含两类问题：对象是否满足条件，以及在满足条件的对象中是否符合个体偏好。价格不超过500元是可以查字段的条件，偏爱简洁外观是需要结合个人信息的倾向。二者可以联合计算，但不能互相抵消：再高的外观偏好分数也不能证明600元的商品满足500元上限。

未知尤其不能被写成满足。假设三个背包价格为450、480、550元，重量分别为0.8千克、未知、0.7千克；用户要求价格不超过500元、重量不超过1千克。第一个通过已知条件，第二个待核验，第三个不通过价格条件。若语义匹配分数依次为0.7、0.9、0.95，分数也没有改变上述状态。搜索会将这类条件进一步联系到查询相关性和资格核验，见第~\ref{chap:rs-search}~章。

偏好匹配则需要把异质证据放到可比较的表示中。标题相近、交互相近和个体喜欢是不同关系，模型通过训练目标学习它们在当前数据中的对应。下面比较协同与文本对齐、知识关系传播、多模态构图和语言内容编码；这些是可替代或组合的建模选择，不是每次请求必须依次经过的步骤。

% Outline: 2.2.3.1
\subsubsection{A-LLMRec：协同与文本表示对齐}
\label{sec:rs-a-llmrec}

协同推荐器知道哪些物品在交互中相关，语言模型却未必能直接读取它的嵌入，两者的向量空间也不互通。\term{A-LLMRec}的主要想法是通过两阶段映射连接两种表示：先对齐协同与文本，再把联合表示投到语言模型的输入空间，用标题生成完成给定协议下的推荐\scite{rs-r23}

\textbf{模型结构。}第一阶段对齐协同与文本：固定预训练协同推荐器，物品协同向量$\bm{v}_i$和文本编码$\bm{t}_i$分别经过编码器$f_I,f_T$，用均方差使$f_I(\bm{v}_i)$接近$f_T(\bm{t}_i)$；各自的解码器重建原向量，避免只靠把所有表示压成同一点降低匹配误差。再用重建后的物品向量与协同用户向量进行正负物品预测。总目标为
\begin{equation}
  \ell_1=\ell_{\mathrm{match}}
      +\alpha\ell_{\mathrm{item-recon}}
      +\beta\ell_{\mathrm{text-recon}}+\ell_{\mathrm{rec}}.
  \label{eq:rs-a-llmrec-stage1}
\end{equation}
原方法在这一阶段微调作为文本编码器的Sentence-BERT；“协同推荐器冻结”不能被扩写成全部语言组件都冻结。第二阶段接入生成语言模型：固定前阶段使用的联合编码器和生成语言模型，训练用户、物品两个投影网络，将表示转换成语言模型词元维度。提示包含用户表示、历史标题及其表示、候选标题及其表示；目标是下一物品标题的逐词元负对数似然。因而更新的是连接网络，不是生成模型全部参数。

已有交互的登山包可走协同编码$f_I(\bm{v}_i)$；全新登山包没有训练过的协同ID向量时，可由文本路径$f_T(\bm{t}_i)$取得相同维度的联合表示。假设两条路径分别输出$(0.8,0.2)$和$(0.75,0.25)$，对齐目标缩小了表示差异，但不保证新包满足重量限制，也不证明两件物品对每个人同样合适。

模型生成标题后，仍需将输出映射到真实目录。重名、别名和目录外文本都要处理；原提示中的候选列表及实验协议也限制了它实际比较哪些物品。标题生成表现不能直接作为完整目录召回能力的证据。

第一阶段的匹配损失把两种编码拉近，重建项分别保持协同与文本信息，推荐项则要求重建后的物品还能与用户因子区分正负项；三者共同防止“向量相同却失去推荐信息”。原实现第一阶段训练10轮，第二阶段5轮，采用Adam；主配置使用SASRec协同表示、Sentence-BERT文本编码和冻结OPT-6.7B。第二阶段只训练到词元空间的投影，梯度可经冻结语言模型传回投影，但不更新生成语言模型本身。服务必须保留协同用户编码器、两条物品联合编码路径、投影和候选提示，不能只部署一个文本嵌入表。

原Amazon实验按用户末次事件测试、倒数第二次验证，每个正例仅配19个随机负物品。表1中Movies and TV的Hit@1为0.6237，SASRec为0.6154、MLP-LLM为0.5838；Video Games却为0.5282，低于SASRec的0.5402。表8删除重建目标使电影Hit@1降至0.5482，冻结Sentence-BERT为0.6173，表9删除联合物品表示降至0.1224，显示对齐路径确有任务作用。论文的“冷物品”主分组是交互量最低35\%而非全部零交互，另有少训练用户和跨目录实验；这几种条件必须分开。其应用证据是小候选商品推荐与标题输出实验，训练或推断速度测量不等同于已经取得生产流量收益\scite{rs-r23}

A-LLMRec靠对齐把协同知识送进生成语言模型，比较范围受提示候选协议限制；当物品本身带有已登记的品牌、主题等显式关系时，KGAT把交互与知识三元组合成一张图，沿关系传播邻居信息。

% Outline: 2.2.3.2
\subsubsection{KGAT：知识与协同关系的注意力匹配}
\label{sec:rs-kgat}

用户只浏览过一个背包，但这个背包的品牌和主题还能连接其他对象，仅凭这一条交互无法利用这些关系。\term{KGAT}的主要想法是把用户—物品交互与物品知识三元组组织成统一图，既学习实体在关系中的匹配，也沿关系传播邻居信息，最终计算用户—物品偏好分数\scite{rs-r84}

\textbf{模型结构。}第一步是知识能量学习：知识三元组$(h,r,t)$表示头实体经关系$r$连接尾实体，关系投影矩阵$\bm{W}_r$与关系向量$\bm{e}_r$给出能量
\begin{equation}
  g(h,r,t)=\lVert\bm{W}_r\bm{e}_h+\bm{e}_r
                         -\bm{W}_r\bm{e}_t\rVert_2^2.
  \label{eq:rs-kgat-energy}
\end{equation}
能量越小，三元组在该表示中越相容，以真实三元组和替换尾实体得到的对照三元组作成对训练。第二步是关系注意传播：邻居传播另用关系相关logit
$a(h,r,t)=(\bm{W}_r\bm{e}_t)^\top
\tanh(\bm{W}_r\bm{e}_h+\bm{e}_r)$，在头实体邻域归一化后加权汇总尾实体。第三步是节点聚合与多层拼接：新节点表示结合自身与邻域，原模型比较相加、拼接和双交互聚合，后者同时变换$\bm{e}_h+\bm{e}_{N(h)}$与$\bm{e}_h\odot\bm{e}_{N(h)}$。多层传播取得多跳关系，拼接各层用户和物品表示后计算内积；训练交替使用知识三元组目标和BPR推荐目标。

在“用户—背包—品牌—另一背包”的小图上，多层传播可让用户表示收到品牌及关联对象的信息。若某节点的两个邻居权重为0.75、0.25，向量为$(1,0)$、$(0,1)$，邻域汇总便为$(0.75,0.25)$，再进入节点聚合，而不是直接成为推荐概率。关系必须已经取得并完成实体链接；错误品牌边也会进入传播。注意权重可以追踪计算依赖，不能被视为某品牌导致购买的因果证明。

知识目标降低真实三元组相对腐化三元组的能量，推荐目标提高已交互物品相对未交互物品的分数；二者共享实体表但约束对象不同。原方法交替用Adam更新关系向量、关系投影、实体嵌入和推荐传播参数，并依据当前关系表示刷新邻域注意。服务保存实体映射、知识与交互边、各层最终表示和评分物品表；图或实体事实改变后，旧表示不会自动体现新边。层数决定传播范围，负尾实体与负物品也须分别采样，不能把两类负例混作一个目标。

原实验在Amazon-book、Last-FM和Yelp2018上用每人80\%随机交互训练、剩余测试，并从训练中另留验证；知识来自Freebase链接或商家属性，评价为全部未训练物品的目录排名。表2中Amazon-book Recall@20为0.1489，NFM为0.1366，NDCG@20为0.1006与0.0913；表5的一层消融删除知识表示或注意均降低指标。层数增加至4的收益趋小，Yelp2018三层还略低于二层，最密集用户组的增益也较弱。应用包括书籍、音乐和本地商家偏好，证据范围是随机交互切分的离线目录匹配，不是时间前推或线上A/B；知识边只能支持其已经登记的关系含义\scite{rs-r84}

KGAT依赖人工登记并完成实体链接的知识三元组；当物品之间没有显式登记关系、却能由图像或文本发现相似性时，LATTICE从模态特征学习物品图再传播ID表示。

% Outline: 2.2.3.3
\subsubsection{LATTICE：多模态物品图的偏好匹配}
\label{sec:rs-lattice}

KGAT的知识边需要人工登记；有些物品关系没有登记，却能由图像或文本特征发现。\term{LATTICE}的主要想法是按各模态特征建立稀疏物品图，学习特征投影以调整邻域，再传播物品ID表示，将结果与协同表示结合用于推荐。内容在这里主要决定边，而不是直接取代全部协同计算\scite{rs-r85}

\textbf{模型结构。}第一步是构图与融合：对于模态$m$，由原特征的余弦相似度构建初始近邻图$\bm{A}^{(m)}_0$；可学习投影产生新的相似图$\widetilde{\bm{A}}^{(m)}$。经过稀疏化和归一化后，混合初始关系与学习关系，再以可学习模态权重融合：
\begin{equation}
  \bm{A}=\sum_m\pi_m
   \bigl[\lambda\bm{A}^{(m)}_0
          +(1-\lambda)\widetilde{\bm{A}}^{(m)}\bigr],
  \qquad \sum_m\pi_m=1.
  \label{eq:rs-lattice-graph}
\end{equation}
第二步是沿图传播ID嵌入并评分：沿融合图传播的是物品ID嵌入，将图分支表示加入协同模型的物品表示，与用户表示形成得分，以BPR目标联合训练模态投影、融合及推荐参数。

假设两款外观相近的背包在视觉图中的边权为0.9，文本用途差异使文本边权仅为0.2。暂不考虑图归一化变化，若视觉、文本权重为0.25、0.75，融合权重为0.375；改成0.75、0.25，则为0.725。这个局部算例说明模态取舍怎样改变传播，并没有证明较大的权重更正确。学习仍需看最终偏好监督是否支持该邻域。

相似图可能把旅行箱外观和通勤背包联系起来，却没有核验用途或互补性。新物品也需要可用内容、图节点及表示更新路径；在已有目录图上评测无交互节点，与无限接收训练后新节点是不同条件。模型输出匹配分数，不能代替属性事实检查。

图像与文本特征由已有内容编码器取得，LATTICE训练的是模态投影、图融合权重、物品ID传播和协同网络，而非把原始像素编码器全部重训。Top-$K$稀疏化改变邻居集合，当前保留边的相似度可接受梯度，但被丢弃边不会因此获得普通连续梯度。初始图提供固定结构参照，学习图随投影调整；上线至少需要版本一致的模态特征、图和最终物品表示，否则内容更新与偏好评分不在同一空间。

原Amazon Clothing、Sports和Baby实验随机按每人80/10/10划分，另移除20\%物品全部训练交互测试冷物品；这些冷节点的内容仍可用于构图。表2中Clothing的Recall@20为0.0710，GRCN为0.0631；Sports为0.0915与0.0833。表3将同样图传播改为传播内容特征，Clothing为0.0643，低于传播ID的0.0710；直接把内容加到LightGCN更低，并在Sports不如原LightGCN。近邻数继续增大也会因混入更多边而下降，说明模态有用与每个模态边都可靠不是同一判断。论文提供商品多模态与无交互节点的离线证据，没有用这些表证明训练后新目录节点已在生产中无缝接入\scite{rs-r85}

LATTICE仍沿物品ID嵌入传播，ID绑定原目录；当迁移到新目录、ID不可复用而属性文字可用时，Recformer直接把键值属性文本作为物品表达。

% Outline: 2.2.3.4
\subsubsection{Recformer：物品文本与历史的联合匹配}
\label{sec:rs-recformer}

LATTICE仍沿物品ID传播，而ID嵌入把物品绑定到某个目录；迁移到新目录时，即使两件商品描述相似，也没有自然共享的ID。\term{Recformer}的主要想法是直接把键值属性文本作为物品表达，再把历史物品文本组成用户输入，以共享语言编码器学习可比较的历史与物品向量\scite{rs-r86}

\textbf{模型结构。}第一步是文本编码：输入不只是把若干标题无分隔地连接起来，词元嵌入之外，类型、词元位置和物品位置嵌入区分属性结构及行为归属；Longformer在局部注意窗口及全局位置间交换信息，得到历史表示$\bm{h}_u$。同一编码结构处理单个物品文本，得到$\bm{v}_i$。第二步是相似度匹配：用温度缩放的余弦相似度预测下一物品。预训练结合掩码语言任务和物品对比任务：前者恢复被隐藏的属性词元，后者将历史表示与真实下一物品拉近、与采样物品区分。推荐微调还要维护目录表示，原方案先随训练阶段更新物品表示表，再固定选定的物品表继续适配历史编码；若编码器已经改变而物品缓存完全未更新，匹配空间可能错位。

例如，历史背包文本包含“品牌：A；材质：尼龙；用途：通勤”，新物品包含“品牌：B；材质：尼龙；用途：通勤”。文本路径让两者能共享材质和用途信号。假设归一化历史向量为$(1,0)$，新包和相机向量分别为$(0.8,0.6)$与$(0,1)$，余弦得分为0.8和0。真实模型从文本和反馈学习这些向量，不是预先把通勤指定成第一个坐标。

服务时先编码当前目录，历史编码后通过相似度访问物品；这与自由生成标题不同。双向编码仅作用于已经发生的历史，不能读入未来目标的文本作为历史事件。文本遗漏关键属性或宣传词压过实际用途时，相似性仍可能失准，需要单独检验内容质量与条件符合情况。

预训练以Longformer参数初始化，共享编码器同时接受掩码属性词元与后继物品对比梯度。目录物品向量先离线编码成表；微调第一阶段每轮按当前编码器刷新这张表，再学习历史与正物品的匹配，选择效果较好的表后固定，第二阶段继续适配历史编码。目录表固定不等于冻结词嵌入或全部历史编码器。原配置单物品每属性最多32词元、历史总计1,024词元，因而50件物品上限不保证每件长描述都完整进入输入；服务要使用相同截断与属性结构。

原文以7个Amazon目录约3,360万交互预训练，另一个目录验证，在6个不同目录微调；所有后继评价均为完整目录而非少量采样候选。表2中Scientific的NDCG@10为0.1027、Recall@10为0.1448，UniSRec为0.0862、0.1255；Instruments的Recformer Recall@10为0.1052，却低于UniSRec的0.1119，不能写成每项皆优。表4目录表维护及预训练消融显示收益随目标目录不同；冷物品实验还对SASRec提供共同冷词元先验，而非纯随机向量。应用是跨商品目录的文本后继匹配与零交互物品访问，原研究未据这些离线结果报告线上业务收益\scite{rs-r86}

Recformer把整段历史的物品文本交给同一个编码器，多人共同浏览同一物品时会重复处理同一段描述；HLLM把物品内容编码与用户行为编码分成两层，让物品向量只算一次、由各段历史复用。

% Outline: 2.2.3.5
\subsubsection{HLLM：语言内容与行为的分层匹配}
\label{sec:rs-hllm}

Recformer把所有历史物品的完整文本交给同一个编码器，多人共同浏览同一物品时会反复处理同一段描述。\term{HLLM}的主要想法是将物品内容编码和用户行为编码分为两层：Item LLM从物品文本及专用词元取得向量，User LLM把这些向量作为序列输入，学习下一物品或候选响应\scite{rs-r87}

\textbf{模型结构。}第一步是物品内容编码：对物品文本$T_i$，第一层给出$\bm{v}_i=f_{\mathrm{item}}(T_i)$，物品向量与用户无关，可以先算一次再由各段历史复用。第二步是行为序列编码：第二层读取请求前的$\bm{v}_{i_1},\ldots,\bm{v}_{i_n}$，形成用户输出$\bm{h}_{u,t}$。下一物品任务用实际后继和随机负例构成InfoNCE目标，即提高正物品在正负集合内的归一化相似度，其形式与式~\eqref{eq:rs-pinnerformer-loss}相近，但监督是后继物品，不是同一个未来窗口任务。训练更新参与适配的两层模型。

候选响应任务有两种接口。早融合把候选向量追加到用户序列，让候选参与深层计算，再接点击判别头；后融合先独立编码历史和候选，在输出侧进行交互。后融合可以复用同一用户表示，物品与用户向量都可预存；早融合需要保留各候选的相关计算，不能复用一份用户表示。二者以点击交叉熵训练，还可结合下一物品辅助目标；不能因共用Item LLM就忽略使用接口差异。

假设100名用户都浏览同一款背包，冻结物品编码器后，其内容向量可以先算一次，再由100段不同历史复用。若下一物品输出为$(0.6,0.8)$，背包与相机的归一化物品向量分别为$(1,0)$、$(0,1)$，相似度为0.6和0.8；这只是候选匹配结果，不是模型生成了一段相机介绍。物品文本改变或编码器继续更新时，缓存也须随版本刷新。

层次编码减少重复内容处理，不保证内容向量保留一切细节，也不能把原HLLM的向量预测扩写成广告创意生成。文本事实、训练时可用历史和目录映射仍须各自核验。

原上线训练分三阶段。先联合更新Item LLM和User LLM，以候选后融合响应目标适配，历史截断150条；再冻结Item LLM并建立目录表，只更新User LLM，把历史扩到1,000条；最后冻结两层，抽取用户与物品特征给现有线上排序器训练。服务在物品创建时编码内容，每天更新此前一天活跃用户的表示，保存两张向量表。它不同于每次请求运行两个大模型，也不同于从未经过推荐适配就冻结物品语言模型；后者在消融中效果很低。

PixelRec与Amazon Books离线使用留一后继评价，HLLM-1B的两层均以TinyLlama-1.1B初始化。表7中Pixel8M的Recall@10为6.129\%，同配方SASRec为5.083\%、HSTU-1B为5.120\%；增大负例数和批次的结果另列，不能把全部差异归于参数规模。表4冻结未适配Item LLM得到0.945\% Recall@10，联合适配为5.581\%；但表8先适配再冻结缓存取得5.218\%，低于全部微调的6.106\%，说明缓存仍有质量取舍。论文确实报告工业排序A/B及部署，称关键指标提高0.705\%，但未充分披露指标名称、周期与置信区间。本节保留部署事实，不把这个数字擅自命名为CTR或销售额\scite{rs-r87}

这些内容模型的共同问题是“已有历史怎样与对象证据比较”，但证据进入位置不同。表~\ref{tab:rs-content-model-tradeoffs}把信息源、学习约束和服务输出放在同一坐标中；高匹配分仍不能替代价格、重量或兼容性事实核验。

\begin{table*}[htbp]
  \centering
  \small
  \begin{tabularx}{\linewidth}{@{}lXXX@{}}
    \toprule
    模型 & 内容进入位置 & 主要学习约束 & 服务使用与边界 \\
    \midrule
    A-LLMRec & 协同与文本联合编码、语言提示 & 对齐、重建、推荐及标题似然 & 候选标题输出；仍需真实目录映射 \\
    KGAT & 实体与关系边 & 知识能量、注意传播及BPR & 用户物品向量；依赖实体链接与图版本 \\
    LATTICE & 模态相似物品图 & 图投影融合、ID传播及BPR & 图与协同向量；相似边不是事实关系 \\
    Recformer & 属性文本与历史词元 & 掩码词元、物品对比、分阶段微调 & 目录相似度；需刷新物品表示表 \\
    HLLM & 独立物品编码、上层行为编码 & 后继对比或候选响应 & 向量或响应；缓存需先适配与版本一致 \\
    \bottomrule
  \end{tabularx}
  \caption{内容与协同匹配的五种接口。输入文字、已链接知识和模态相似图提供不同依据，不能互换为已经核验的对象属性。}
  \label{tab:rs-content-model-tradeoffs}
\end{table*}

已有可靠实体关系时，KGAT能把关系含义带入多跳传播；仅有图片和文本时，LATTICE学习的是相似邻域，需要另外检查用途与类别混淆。新目录ID不可复用但属性文字可用时，Recformer或HLLM的内容路径能提供表示，代价是内容编码及目录表维护。要把成熟协同知识接入生成语言模型，A-LLMRec的对齐与提示更直接，但小候选标题评价不能替代大目录访问。选用内容路径还要检查信息的新鲜度、缺失字段及广告宣传文本，不能只依据语言模型规模判断可靠性。

% Outline: 2.3
\section{信息获取与维护}
\label{sec:rs-evidence-maintenance}

已有信息足够时，问题是怎样编码；缺少或过时的信息则需要先处理证据本身。补充提供新的依据，更新让依据反映当前状态，检验确定依据能支持什么结论。这三项职责贯穿用户、对象和匹配建模，不能等评分器训练完才统一补做。

% Outline: 2.3.1
\subsection{缺失信息补充}
\label{sec:rs-missing-information}

“\term{冷启动}”常把不同缺口混在一起。新用户可能没有历史却愿意填写用途；匿名用户可能有当前会话却没有稳定身份；新物品可能没有交互却有完整属性；长期低曝光物品可能存在很久，仍缺少可比较的反馈。训练时未见实体又是模型参数覆盖条件，和有没有内容不是同一回事。先指出缺什么，才能选择内容先验、显式初始化、关系共享或少量反馈适应。

已有其他业务的信息时，还需建立对应关系。跨域推荐涉及不同目录或交互分布，至少要说明用户或物品是否重叠、迁移单向还是双向、目标域零反馈还是少反馈。同一商品目录在首页和购物车的CTR适配是多场景问题，留到第~\ref{chap:rs-ranking}~章讨论。不能因为都叫“迁移”，就默认身份、对象和行为含义已经对齐。

能询问用户时，补充信息也可以直接来自回答。例如，不清楚“轻便”指重量上限还是肩带舒适度，询问后可以改变可选集合；若两个解释最终产生完全相同的选择，继续询问的收益可能很小。问题价值要和打断成本一起衡量，回答还需带上适用范围。用户此次答“买给同事”，不应被永久登记为本人偏好改变。交互界面与展示测试分别由第~\ref{chap:rs-presentation}~章承接。

% Outline: 2.3.1.1
\subsubsection{DropoutNet：偏好信息缺失下的内容补充}
\label{sec:rs-dropoutnet}

已有潜在模型在热物品上有效，但新物品没有学好的交互向量，直接查表只能得到随机或零表示。\term{DropoutNet}的主要想法是从这个教师模型出发，把用户和物品的偏好向量与内容特征分别送入映射网络；训练时随机移除整侧偏好输入，让网络学习在该输入缺失时依靠内容恢复原模型的相关性得分\scite{rs-r02}

\textbf{模型结构。}第一步是双侧映射：设原向量为$\bm{p}_u,\bm{q}_i$，内容为$\bm{c}_u,\bm{c}_i$，随机掩码$m_u,m_i\in\{0,1\}$控制偏好侧输入，新表示及恢复目标可写成
\begin{align}
  \widehat{\bm{p}}_u&=f_U([m_u\bm{p}_u;\bm{c}_u]),&
  \widehat{\bm{q}}_i&=f_I([m_i\bm{q}_i;\bm{c}_i]),\\
  \ell_{u,i}&=\bigl(
    \langle\widehat{\bm{p}}_u,\widehat{\bm{q}}_i\rangle
      -\langle\bm{p}_u,\bm{q}_i\rangle\bigr)^2.
  \label{eq:rs-dropoutnet}
\end{align}
第二步是缺失模拟：有偏好输入的样本促使网络保留原模型的信息，偏好被移除的样本迫使内容参与恢复。丢弃的是整个偏好向量，不等于任意隐藏层的常规dropout。

假设原模型给一位用户和旧登山包的得分为0.8，训练将包的偏好向量置零，内容网络给出的新内积为0.6，则平方误差为0.04，梯度推动内容路径解释缺失的部分。新登山包服务时同样以零偏好向量和真实属性输入。新用户取得少量交互后，原方案还以交互过物品的表示构造偏好输入近似，并在训练中模拟这种早期状态。

输出是新向量与匹配分数，能够接入内积召回；其监督目标仍来自原潜在模型。原模型有偏差、内容无信息或内容与交互的对应关系改变时，恢复教师得分也不保证恢复真实偏好。两侧都缺少可用信息时，输入缺失模拟不能凭空补出依据。

原潜在模型先训练并固定，然后映射网络以其得分为监督。每个批次随机抽取用户或物品，交替把偏好输入置零或替换为已交互对象因子的均值，使零反馈与初期反馈两条服务路径都在训练中出现。均值只是暂时近似，不是重新求解原潜在模型；原模型重训后应重新使用匹配版本的因子。服务保留两侧映射、内容、教师因子及反馈集合，没有内容的一侧不能仅靠置零恢复原ID知识。

CiteULike实验含5,551用户、16,980论文和约20.5万收藏，内容是标题摘要的8,000词特征，移除3,396篇论文训练交互作冷物品。表1中以WMF为教师的DropoutNet冷启动Recall@100为0.636，DeepMusic为0.601、CDL为0.573；热启动为0.593，低于CDL的0.603。丢弃率增加到0.7以上还导致热启动明显下降。XING职位数据另测试用户与职位冷启动，深层网络提高冷启动却略损热启动，揭示内容补充与保留原协同效果的取舍。它们是论文和职位推荐的真实数据离线实验，不是已经观测到招聘转化或线上收入改善\scite{rs-r02}

DropoutNet只利用已经登记的内容特征恢复教师得分；当物品资料与行为记录仍遗漏领域中的选择因素时，KAR借助语言模型离线生成偏好推理与物品知识文本，再适配为推荐特征。

% Outline: 2.3.1.2
\subsubsection{KAR：语言知识与行为的适应匹配}
\label{sec:rs-kar}

DropoutNet只恢复已有内容特征，而商品资料可能未覆盖领域中的全部选择因素，行为记录也未直接说明用户为何选择。\term{KAR}的主要想法是按领域因素组织提示，让语言模型离线生成用户偏好推理文本和物品知识文本，再将这些文本编码并适配为推荐特征\scite{rs-r88}

\textbf{模型结构。}第一步是离线生成与编码：偏好提示以已有历史为条件，知识提示以物品信息及领域因素为条件，离线文本编码得到$\bm{t}_u,\bm{t}_i$。第二步是混合专家适配：适配器包含共享专家和分别服务用户推理、物品知识的专属专家，用户侧可概括为
\[
  \bm{z}_u=\sum_{e\in E_{\mathrm{shared}}\cup E_{\mathrm{user}}}
                \alpha_e(\bm{t}_u)f_e(\bm{t}_u),
\]
物品侧同理但使用自己的专属专家。门控权重$\alpha_e$选择各专家贡献；适配输出与既有用户、物品及情境字段一起进入推荐网络，随点击等下游目标训练。语言生成和编码可预先完成，服务时读取特征。

假设背包知识补充了容量、面料和使用场景三个方面，适配器中共享与专属专家的教学输出分别为$(0.6,0.2)$和$(0.2,0.8)$，门权重为0.25和0.75，则输出为$(0.3,0.65)$。这个向量只是评分的新输入，尚未作出推荐。若生成文本错误地把普通防泼水写成可完全浸水，预测网络也不会自动把它纠正成真实规格。

因此，物品事实需要核验来源，偏好推理应保留为待检验假设。丰富的描述和流畅的理由不构成用户真实动机的证据，更不构成因果解释。相较第~\ref{sec:rs-a-llmrec}~节的表示对齐，KAR主要增加外部生成内容，二者信息来源及验证责任不同。

生成语言模型与知识编码器先离线调用，点击或重排梯度只训练混合专家适配器与下游网络；不会凭一次点击自动纠正文中错误事实。原专家是两层MLP，门网络按输入分配共享和专属专家权重，以便用户推理和物品描述既共享某些语义又保持职责差异。适配后的低维向量还可预存；此时服务不执行生成或专家网络，只读版本对应的特征。未生成知识的新用户或物品先用默认向量，后续离线补齐，不能把默认值当作已经取得知识。

MovieLens-1M实验把4至5分作为正标签、其余负标签，按用户ID划分；Amazon Books约140.7万交互将5分作为正标签，二者不是曝光点击。表1中DIN加KAR在Books的AUC从0.8304升至0.8418，LogLoss从0.4937降至0.4801；用户推理和物品事实单独及联合的消融说明两类输入各有贡献。与图知识一起使用并未持续超过只用语言知识，不能据此称二者必然互补。真实部署在Huawei新闻和音乐：音乐7天A/B每组10\%用户，歌曲播放量提高1.7\%、播放总时长提高1.57\%；新闻另报告召回指标提高7\%，但未充分定义该召回口径。部署证据说明生成特征可接业务，不认证每条生成理由是真实用户动机\scite{rs-r88}

KAR补充的是外部生成的内容先验，偏好判断仍要落到真实反馈；新用户只有几次评分、不足以独立训练完整网络时，MeLU从已有用户学习一份少量反馈即可适应的初始化。

% Outline: 2.3.1.3
\subsubsection{MeLU：少量反馈的用户偏好适应}
\label{sec:rs-melu}

KAR补充内容先验，但偏好判断仍要落到反馈；新用户有了几次真实反馈，仍不足以为其独立训练完整网络。\term{MeLU}的主要想法是从已有用户学习一份容易适应的初始化：把每个用户作为一项训练任务，用少量支持交互调整预测层，再用其他交互检查调整是否有用\scite{rs-r89}

\textbf{模型结构。}第一步是内层局部适应：把参数分为属性编码部分$\theta_{\mathrm{emb}}$和偏好预测部分$\theta_{\mathrm{pred}}$，对用户$u$的支持集$S_u$进行局部更新：
\begin{equation}
  \theta'_{\mathrm{pred},u}=
  \theta_{\mathrm{pred}}-\alpha\nabla_{\theta_{\mathrm{pred}}}
     \ell(S_u;\theta_{\mathrm{emb}},\theta_{\mathrm{pred}}).
  \label{eq:rs-melu-inner}
\end{equation}
此时保持属性嵌入部分不变。第二步是外层元优化：随后在该用户查询集$Q_u$上评价适应后的预测，跨用户最小化
$\sum_u\ell(Q_u;\theta_{\mathrm{emb}},\theta'_{\mathrm{pred},u})$，
更新共享初始化及属性编码。支持集与查询集承担不同训练职责，不能用同一反馈反复证明适应后泛化更好。

用标量预测层说明一步适应。假设背包特征$x=1$、评分$y=4$、初始参数$w=3$，损失为$\tfrac12(wx-y)^2$，步长为0.2。梯度为$-1$，更新后$w'=3.2$；特征$x=0.5$的候选预测从1.5变为1.6。真实MeLU以用户和物品属性进入多层预测器，这里仅隔离局部梯度怎样改变后续判断。

服务时根据新用户少量标注执行局部更新，再对候选评分。完全没有反馈时，式~\eqref{eq:rs-melu-inner}无从执行，只能使用共享先验。元训练用户与新用户差异很大、支持反馈噪声较多时，快速适应也可能快速过拟合；取得哪些反馈仍是独立的信息获取问题。

原任务是显式评分回归，以支持集评分误差执行1至5次局部更新，再让查询集评分误差经局部更新路径反传到共享初始化。属性编码在内层冻结，外层仍可更新；用户特有的是适应后的预测参数，不是另建一张永久用户ID表。服务需要保留初始化、属性映射与该用户支持标签，支持集变化后重新适应或按约定增量更新；训练查询集只用于外层评价，不能成为新用户当次推断输入。

MovieLens与Bookcrossing原实验按电影或图书年份分新旧对象，再随机分新旧用户，每分区只保留13至100条评分用户，10条作查询、余下作支持。表2中MovieLens新用户旧物品的MeLU-1平均绝对误差为0.7884，Wide \& Deep为1.0694；5次更新为0.7854，改进很小。Bookcrossing的相同条件下1次更新1.8701，5次反而1.8767，说明多更新不必更好。论文还做84名用户的电影询问演示，比较选择信息性证据与热门电影，并报告后续推荐改善；这是真实小规模用户研究，不是生产持续A/B。后续复核强调必须给非元学习基线同样的支持反馈及调参机会，元初始化的价值不能由不公平输入比较决定\scite{rs-r89,rs-ch02-cold-meta-review}

MeLU在用户侧适应预测层，新广告的问题在对象侧：它有图片、类别和属性，却没有成熟的ID向量。MetaEmb学习由内容生成初始ID嵌入的生成器，适应对象与更新参数的规模都与MeLU不同。

% Outline: 2.3.1.4
\subsubsection{MetaEmb：新广告嵌入的内容初始化与反馈适应}
\label{sec:rs-metaemb}

MeLU在用户侧适应预测层；新广告的问题在对象侧：它已经有图片、类别和商品属性，却没有成熟的ID向量。\term{MetaEmb}的主要想法是学习一个由内容产生初始嵌入的生成器，并同时考察它在零反馈时的预测和少量反馈更新后的预测。与MeLU主要适应用户预测层不同，这里适应的是新广告的ID嵌入\scite{rs-a28}

\textbf{模型结构。}第一步是内容初始化：固定已有CTR主网络参数$\theta$，生成器$g_\phi$由广告特征$\bm{c}_i$给出$\bm{v}^{(0)}_i=g_\phi(\bm{c}_i)$。第二步是内层反馈更新与外层双目标：在旧广告上模拟新ID，取两组少量样本$D_i^a,D_i^b$，计算
\begin{align}
  \bm{v}^{(1)}_i&=\bm{v}^{(0)}_i
      -\eta\nabla_{\bm{v}^{(0)}_i}\ell(D_i^a;\theta,\bm{v}^{(0)}_i),\\
  \ell_{\mathrm{meta},i}&=
     \lambda\ell(D_i^a;\theta,\bm{v}^{(0)}_i)
      +(1-\lambda)\ell(D_i^b;\theta,\bm{v}^{(1)}_i).
  \label{eq:rs-metaemb}
\end{align}
生成器通过两项损失学习：初始向量应立即可用，也应在少量样本更新后表现合理。主CTR网络在这一训练过程中冻结，不能把其变化误记为初始化效果。

假设某新背包广告的初始标量嵌入为0.2，教学CTR头为$\sigma(v)$。收到一次点击后，交叉熵对$v$的梯度为$\sigma(0.2)-1\approx-0.450$；步长0.1时，更新到约0.245，预测由0.550变为约0.561。真实投放中的其他字段和多条样本会改变这个结果，算例仅区分内容初始化与反馈适应两步。

没有反馈时只执行初始化，反馈到来后才更新ID向量。推广到普通新物品时需要明确CTR任务和字段的对应；广告资格、支付和预算都不由这个嵌入生成器决定。内容缺失及新广告分布变化也会限制初始化的可靠性。

原生成器复用成熟CTR模型的广告属性嵌入，冻结这些表，对字段向量池化后用一个带tanh、无偏置的稠密层生成新ID。外层训练使生成器在初始和内层一步更新后的两份样本上都可用，需经内层梯度传播，不能只训练成模仿旧ID的回归器。上线新ID先由内容生成，真实反馈到来后只更新ID，主网络保持本阶段的固定版本；因此本方法适应的参数规模远小于MeLU的预测层。

原实验包含MovieLens二值评分、Tencent应用广告CVR和KDD Cup 2012搜索广告CTR。它按记录量将旧广告与新广告分开，并给每个新广告连续3批反馈；Tencent与KDD每批200条，不能概括为仅有一次点击。表2汇总6个主模型、每个3次运行，Tencent冷启动AUC相对随机ID基线提高0.50\%，经过3批后随机ID和MetaEmb相对同一初始基线分别提高1.57\%、1.83\%；后两个数字不是直接的1.83\%相对优势。MovieLens增益更大，实际广告日志较小，且没有完整上线A/B报告。其应用证明是新广告初始化与早期反馈学习，不能写成已验证预算或投放收益\scite{rs-a28}

MetaEmb在同一目录内由内容初始化新ID；当偏好证据来自另一个目录或业务、目标目录本身缺少记录时，EMCDR先分别学习两域因子，再借重叠实体训练跨域映射。

% Outline: 2.3.1.5
\subsubsection{EMCDR：跨域表示映射的偏好补充}
\label{sec:rs-emcdr}

MetaEmb在同一目录内初始化；当同一用户在来源目录有充足交互、在目标目录却没有记录时，可以尝试把已知偏好映到目标空间。\term{EMCDR}的主要想法是分别在两域用矩阵分解或BPR学习潜在因子，再用重叠用户或物品训练映射。两域因子学习和跨域映射是两个阶段\scite{rs-r95}

\textbf{模型结构。}第一步是两域独立因子学习：来源域与目标域分别拟合MF评分误差或BPR正负次序，不要求潜在坐标直接可比。第二步是重叠实体映射：以重叠用户集合$U_{\cap}$为例，来源和目标因子为$\bm{p}_u^S,\bm{p}_u^T$，映射$f_\phi$通过
\begin{equation}
  \min_\phi\sum_{u\in U_{\cap}}
     \lVert f_\phi(\bm{p}_u^S)-\bm{p}_u^T\rVert_2^2
  \label{eq:rs-emcdr-map}
\end{equation}
学习，$f_\phi$可以是线性映射或MLP。第三步是目标侧评分：对仅在目标侧缺交互的用户，以$f_\phi(\bm{p}_u^S)$和目标域已有物品因子计算内积，再进入目标目录的Top-$K$访问。重叠物品可建立相应的物品映射任务，但不是所有应用都会同时具有两种重叠。

假设来源用户因子为$(1,0)$，线性映射为$\operatorname{diag}(2,1)$，目标用户表示为$(2,0)$。目标候选向量为$(1,1)$与$(0,2)$，得分分别为2和0。映射改变的是表示坐标及关联，不表示来源域的评分2与目标域评分2拥有相同原始含义。

这个过程需要来源因子、对齐实体和可学习的跨域关系。没有重叠实体，式~\eqref{eq:rs-emcdr-map}就缺少监督；两域偏好关系变化时，拟合良好的旧映射也可能造成负迁移。只能取得聚合数据或身份使用受限时，需要另建可用证据条件，不能默认原方法仍可直接应用。

两域先各自拟合MF评分误差或BPR正负次序，不要求潜在坐标直接可比；映射阶段固定这些因子，以两域均有足够反馈的对齐实体训练MLP或线性映射。选择反馈充分的映射实体可降低目标因子的噪声，但也可能偏向活跃用户，迁移到稀疏用户仍需检验。服务保留来源因子、映射与目标物品表；目标新物品若也没有因子，就不能只映用户完成匹配。原映射MLP用输入输出同维、隐藏宽度两倍的结构，以小批量梯度迭代求解。

MovieLens到Netflix实验是约5,871部重叠电影而非共享用户；Douban电影到图书是10,000个共享用户而非相同目录。原文在目标域随机移除10\%至50\%实体全部评分，各重复10次，MF用RMSE、BPR用AUC分别评价。表3在20维、10\%电影冷启动时非线性MF映射RMSE为0.8988，线性映射0.9160、CMF为0.9251；表4非线性BPR映射AUC为0.6892、线性0.6730。随着移除比例增长效果降低，两域因子质量与对齐量共同限制迁移。证据是电影、图书评分的跨域离线补充，既不覆盖无重叠实体，也未测真实曝光CTR\scite{rs-r95}

EMCDR只在训练时借重叠实体学一份静态映射，运行时两域表示并不交换；CoNet在网络层直接跨域连接，让一域的隐藏表示参与另一域计算，共享方式与静态映射不同。

% Outline: 2.3.1.6
\subsubsection{CoNet：跨域网络连接的偏好共享}
\label{sec:rs-conet}

EMCDR只在最终因子处用一份静态映射交换信息，中间层的域内表示并不互通。\term{CoNet}的主要想法是让两个目录的预测网络在多层计算中共享隐藏信息。原设定以相同用户连接两域，各域保留独立物品表和预测头，通过跨域连接联合训练\scite{rs-r96}

\textbf{模型结构。}第一步是双分支并行计算：对两个域$S,T$，一对连接层可写为
\begin{align}
  \bm{h}^{S,(\ell+1)}
    &=\rho(\bm{W}_S^{(\ell)}\bm{h}^{S,(\ell)}
           +\bm{H}_{T\to S}^{(\ell)}\bm{h}^{T,(\ell)}
           +\bm{b}_S^{(\ell)}),\\
  \bm{h}^{T,(\ell+1)}
    &=\rho(\bm{W}_T^{(\ell)}\bm{h}^{T,(\ell)}
           +\bm{H}_{S\to T}^{(\ell)}\bm{h}^{S,(\ell)}
           +\bm{b}_T^{(\ell)}).
  \label{eq:rs-conet-cross}
\end{align}
$\rho$为激活函数；$\bm{H}$传递另一域当前层的信息，$\bm{W}$处理本域信息。第二步是跨域连接与联合训练：原CoNet令两个方向共享同一矩阵，即$\bm{H}_{T\to S}^{(\ell)}=\bm{H}_{S\to T}^{(\ell)}$，这里的方向下标只是区分两次调用。两域的真实交互及采样未观测物品提供各自的二元交叉熵，共同更新共享用户因子、域内网络与连接；SCoNet再对连接矩阵施加稀疏约束，以减少不必要的信息传递。

假设本域变换和偏置合计0.3，跨域贡献为0.2，则ReLU输出0.5；把连接置零，输出回到0.3。这可作为局部消融的教学计算，实际负迁移应比较目标域完整预测，而不能以非零连接数量判断。服务时还需按训练定义提供跨域分支的可用输入，再由目标域预测头比较物品；共享一个用户因子并没有消除候选相关网络。

模型需要原设定中的用户对应和两域训练信息。连接稀疏不保证迁移方向正确，也不覆盖任意无重叠目录或目标域零反馈任务。若只有来源历史而无法构造训练时使用的另一域输入，应重新设计接口并评价，不能把缺失分支静默当作已经共享的知识。

原训练为两个域分别组织同用户的正交互及采样负物品，按域轮流更新各自主网络，共享用户表与跨域连接受两域梯度约束。SCoNet加入连接矩阵的$\ell_1$惩罚，使部分传递权重收缩到零；它改变传递哪些隐藏坐标，而不是删除全部来源历史。原论文没有提供适用于任意线上请求的另一域候选构造协议，落地时必须规定配对输入或改成明确的来源历史汇总，并按新接口重训；直接将缺失向量填零是新的服务选择，不是原实验已经验证的默认保证。

原Mobile数据以23,111共享用户的新闻阅读辅助应用安装，Amazon以80,763共享用户的电影评分辅助图书偏好，留一目标交互并与99个负物品排名。表2中Amazon的SCoNet NDCG@10为0.3424，CoNet为0.3261、仅共享用户的MLP++为0.3175；但CMF为0.2378，低于不迁移MLP的0.3143，直接展示负迁移。Mobile虽SCoNet数值更高，对最佳基线的配对检验$p$值为0.20至0.34，不能称统计显著。表3少量减少目标反馈后SCoNet也会降到原单域基线以下。应用证据因此支持有条件的跨域共享，不证明稀疏矩阵能普遍避免负迁移，也没有给出生产请求完整性能报告\scite{rs-r96}

从DropoutNet的内容恢复、KAR的生成先验、MeLU与MetaEmb的元适应，到EMCDR的静态映射与CoNet的层间连接，六者补充的缺口位置不同：缺内容、缺反馈、缺对象还是缺跨域证据，决定了各自能更新什么、零反馈时能否启动。

信息缺失方法应按缺口选择。表~\ref{tab:rs-information-completion}列出它们必须额外取得什么、具体更新什么，以及零反馈是否已经可用。所有方法都依赖先验或其他训练数据，“零反馈可用”不等于“没有证据也知道偏好”。

\begin{table*}[htbp]
  \centering
  \small
  \begin{tabularx}{\linewidth}{@{}lXXX@{}}
    \toprule
    模型 & 必需补充信息 & 训练及适应对象 & 使用条件 \\
    \midrule
    DropoutNet & 成熟偏好模型与有信息的内容 & 两侧映射；模拟偏好输入缺失 & 零反馈走内容，少反馈用因子均值近似 \\
    KAR & 可用描述、历史与生成知识 & 专家适配及下游网络 & 生成可离线，事实与推理需核验 \\
    MeLU & 用户属性与少量评分 & 外层初始化，内层用户预测层 & 无评分不能执行局部适应 \\
    MetaEmb & 广告属性与成熟主网络 & 生成器，随后新ID向量 & 零反馈先初始化，反馈后更新小向量 \\
    EMCDR & 来源因子及可靠重叠实体 & 独立域因子、随后表示映射 & 目标零反馈可映射，仍需目标对象因子 \\
    CoNet & 对齐用户及两域交互 & 联合主网络及稀疏连接 & 两域输入可用，需检验负迁移 \\
    \bottomrule
  \end{tabularx}
  \caption{信息补充方法的依据与参数更新范围。内容、支持标签和跨域对应是不同资源，不能统称为同一种冷启动设置。}
  \label{tab:rs-information-completion}
\end{table*}

新物品有属性却无交互时，内容初始化或文本表示可以先给出分数；新用户已有少量评分时，MeLU适应预测层，MetaEmb则主要服务少反馈的新广告ID，两者更新对象不可互换。来源目录有记录而目标目录无记录时，EMCDR需要对齐实体，CoNet还依赖联合任务输入。生成知识能够增加描述维度，却必须把事实核验和模型推断分开。比较时应给各基线同样的内容、支持标签和实体对应，否则优势可能只来自多拿了一份信息。

% Outline: 2.3.2
\subsection{理解依据更新}
\label{sec:rs-evidence-refresh}

信息缺失与信息过时需要不同处理。用户刚从登山改为通勤需求，系统可能已经收到明确查询却仍使用旧画像；也可能查询已进入输入，但模型仍过度依赖旧历史。前者是证据接入延迟，后者是计算或学习对变化适应不足，增加更多旧记录未必有帮助。

固定请求发生在10:05，假设用户10:04收藏相机，长期向量上次更新在午夜，模型参数更新在9:00。TransAct若在10:05读到该收藏，可以改变即时分支；长期向量仍未更新并不妨碍这一点，参数也不必在一分钟内重新训练。相反，即使9:00的模型比昨日模型更准确，它也不能使用尚未到达的10:04事件。

MIMN需要检查行为是否已写入记忆，SIM需要检查事件是否已进入可检索历史，TWIN还需检查物品投影与检索参数版本。四种状态不能统一成一个“更新时间”。图~\ref{fig:rs-history-access}比较事件保存和请求读取的差异；信息路径改变之后，延迟、存储及重算成本也会改变。这里借用第~\ref{chap:rs-ranking}~章候选相关历史模型说明时效问题，各模型的完整评分机制不在本节展开。

\begin{figure*}[htbp]
  \centering
  % 三条历史访问路径；最终宽度约162mm，按保存粒度比较。
\begin{tikzpicture}[
  font=\figurefont\footnotesize,
  rsbox/.style={rectangle,draw=BookRule,fill=BookPaper,
    text=BookInk,align=center,text width=40mm,minimum height=13mm,inner sep=4pt},
  rsflow/.style={-{Stealth[length=2mm]},draw=BookMuted,line width=.9pt}
]
  \node[font=\figurefont\footnotesize\bfseries] at (0,1.2) {新事件到来};
  \node[font=\figurefont\footnotesize\bfseries] at (5.7,1.2) {保留的信息};
  \node[font=\figurefont\footnotesize\bfseries] at (11.4,1.2) {预测请求读取};
  \node[rsbox,fill=FigureInput!10] (a1) at (0,0) {增量读写\\擦除、添加、状态更新};
  \node[rsbox,fill=FigureModel!10] (b1) at (5.7,0) {固定容量记忆\\部分事件细节被压缩};
  \node[rsbox,fill=FigureOutput!10] (c1) at (11.4,0) {读取记忆及归纳状态\\结合候选预测};
  \node[rsbox,fill=FigureInput!10] (a2) at (0,-1.9) {追加原始事件\\更新历史索引};
  \node[rsbox,fill=FigureModel!10] (b2) at (5.7,-1.9) {可检索事件记录\\保留逐条行为信息};
  \node[rsbox,fill=FigureOutput!10] (c2) at (11.4,-1.9) {候选检索相关事件\\再作精细聚合};
  \node[rsbox,fill=FigureInput!10] (a3) at (0,-3.8) {更新哈希桶或行为簇\\刷新聚合及版本};
  \node[rsbox,fill=FigureModel!10] (b3) at (5.7,-3.8) {桶或簇表示\\保留聚合信息与数量};
  \node[rsbox,fill=FigureOutput!10] (c3) at (11.4,-3.8) {候选访问桶或簇\\形成聚合兴趣};
  \foreach \r in {1,2,3} {
    \draw[rsflow] (a\r.east) -- (b\r.west);
    \draw[rsflow] (b\r.east) -- (c\r.west);
  }
\end{tikzpicture}

  \caption{长历史的三种访问路径。增量记忆在新事件到来时压缩状态；历史检索在请求时按候选选事件；桶或簇聚合预先改变历史粒度，再由请求取得相关聚合。图比较信息保存与访问职责，不表示各模型的全部预测网络相同。}
  \label{fig:rs-history-access}
\end{figure*}

更新评价应固定请求时点，记录各信息真正可用的时间，比较变化发生后的预测质量和维护代价。目录可用性更新由第~\ref{chap:rs-retrieval}~章承接，模型重训与多场景适应由第~\ref{chap:rs-ranking}~章展开，具体状态同步属于系统卷。更及时的信息可以帮助当前判断，但是否改善长期用户收益仍需后续决策和评价证据。

% Outline: 2.3.3
\subsection{理解质量检验}
\label{sec:rs-understanding-quality}

理解质量不能由画像字数或历史长度证明。用户侧应检查判断是否符合明确表达、近期变化及真实反馈；对象侧应检查属性和关系是否有来源、是否与当前版本一致；匹配侧应检查条件符合情况和偏好预测是否各自成立。生成画像尤其要能回到原始证据，不能让生成内容在后续链条中变成无来源的事实。

比较模型时，先固定证据再改变处理方式。例如，在相同历史上移除顺序信息，检验顺序的贡献；固定事件后改变会话边界，检验划分规则；固定完整历史后改变检索条数，检验选择损失。再单独改变历史长度或信息时效，观察额外证据的贡献。若一次同时增加历史、内容和计算量，最终改善无法归到某一种编码机制。

会话模型的一项独立比较说明为什么要保留简单基线。Ludewig与Jannach在RSC15、TMALL、RetailRocket和Zalando上比较转移、会话近邻和GRU4Rec，表3中RetailRocket顺序会话近邻s-sknn的MRR@20为0.3450，GRU4Rec为0.2427；RSC15却由GRU4Rec取得0.6827的HR@20，高于普通会话近邻0.6213。除一个小数据版本外，结果为5个时间切片平均，GRU层宽限制100、其他参数随机搜索100次，近邻用相同验证集调邻居与样本数。数据、指标与调参预算都限定了结论，不能由其中一行断言神经或非神经方法普遍获胜\scite{rs-ch02-session-evaluation}

序列模型的第二项比较进一步把结构与训练配方分开。Klenitskiy与Vasilev用相同PyTorch训练框架，在5个数据集做全目录留一评价，将SASRec从独立正负二元损失改为目录或3,000负例的分类交叉熵。表2中MovieLens-1M的NDCG@10从0.1341升至0.1857，BERT4Rec为0.1537；Steam则BERT4Rec为0.0662，高于改后SASRec的0.0652。各模型有验证早停并恢复最佳权重，BERT4Rec允许更长训练，不能把未收敛差异当作架构差异。结果说明目标、负例数、收敛和候选评价可以改变模型次序，并没有证实某种注意方向在全部数据上最好\scite{rs-ch02-gold-loss}

表~\ref{tab:rs-understanding-evidence-levels}区分本章实际取得的效果证据。离线排名检验的是日志目标可否被区分，线上A/B测的是完整系统在特定流量中的响应，部署报告说明系统被使用及如何维护；三者各有价值，却不能互相代替。

\begin{table*}[htbp]
  \centering
  \small
  \begin{tabularx}{\linewidth}{@{}lXXX@{}}
    \toprule
    证据层次 & 本章实例 & 实际回答的问题 & 不足以推出 \\
    \midrule
    离线目录排名 & SASRec、BERT4Rec、Recformer、GroupIM & 给定切分和候选集合能否命中目标 & 在线曝光CTR、长期满意或因果动机 \\
    离线响应与消融 & DIN、DSIN、TWIN、MetaEmb & 给定日志标签及预算哪个部件有效 & 收益全部来自架构、任意数据均优 \\
    线上随机对照 & DIEN、SIM、TransAct、SDIM、KAR & 完整部署在指定流量和周期的指标变化 & 单模块净效应、未测人群与长期后果 \\
    部署与维护 & MIMN、TWIN-V2、HSTU、HLLM、LONGER & 哪些状态和计算已接生产、如何复用 & 未披露业务指标的具体含义或普遍收益 \\
    \bottomrule
  \end{tabularx}
  \caption{模型效果的证据范围。原论文AUC、GAUC、HIT与业务指标定义分别保留，不将不同口径的数字排成总榜。}
  \label{tab:rs-understanding-evidence-levels}
\end{table*}

分组检查应覆盖新老用户、稀疏与密集交互、稳定与匿名身份、不同证据来源及变化后的阶段。没有已知事实时，正确输出“未知”可能比给出自信但错误的属性更有用；取得补充信息的成本也要计入。成员注意、图注意和兴趣状态可用于检查网络依赖了什么，但不能独立验证心理或因果解释。

本章检验的是当前依据是否足够、是否一致及是否及时。训练标签如何校正曝光、位置和选择偏差，由第~\ref{chap:rs-ranking}~章建立；完整数据切分、指标与效果识别由第~\ref{chap:rs-evaluation}~章讨论。局部模型比较应保留这些入口，不能以一次离线命中提升结束全部验证。

% Outline: 2.4
\section{本章小结}
\label{sec:rs-understanding-summary}

从登山历史推断通勤请求，关键是分清已知表达、可见行为、对象事实和模型推断。直接条件规定本次需求的一部分，历史和情境补充尚未表达的部分；对象表示提供属性与关系，匹配模型再学习这些证据如何共同支持选择。语义相似和高偏好分数都不能替代硬条件核验，未知信息也应保留为未知。

行为模型在顺序、会话、候选参与、记忆容量和监督时间范围上作出不同选择。它们形成的中间表示及最终输出决定了下一环节怎样使用：候选无关向量适合共享编码，候选相关聚合需要保留逐候选计算，群组请求还需成员关系和共同标签。内容补充、少样本适应及跨域共享各有不同的信息前提，更新与质量检验则限定这些依据在当前请求中是否有效。

原研究的经验结果也须放回这些接口。扩大历史可以增加依据，但记忆槽、检索选集、哈希桶、簇代表和查询压缩各会丢掉不同细节；缓存还把状态维护带到模型之外。内容和元学习在缺失条件下提供可计算的先验，却不消除内容质量、支持反馈与域对应的要求。简单基线、多目标与负例配方都可能改变离线次序，线上收益又来自完整部署。因此，合理取舍是先确认当前任务可用信息与允许损失，再在共同标签、候选和预算下比较，不以论文最大历史长度或单项收益替代选择依据。

现在已有了能够比较需求与对象的表示，但表示本身尚未从大目录中取得真实候选。下一章将把内积、邻域关系与生成输出变成具体访问方法，并检查目录覆盖、资格和候选遗漏。理解形成的证据只有进入这些实际选择，才能继续接受任务效果的检验。
